% Leak-Free but Not Loss-Free: Verifying Permission-Aware Retrieval for RAG
% Venue style is applied later (NeurIPS 2027 E&D style once released); this preprint layout uses the
% standard article class. \tbd{} marks what is still open: the X8 amendment 2 results in e2e.tex, and
% in limits.tex the filing of vendor reports, the human check of references and the archive DOI.
\documentclass[11pt]{article}
\usepackage[margin=1in]{geometry}
\usepackage[T1]{fontenc}
\usepackage{lmodern}
\usepackage{microtype}
\usepackage{amsmath,amssymb,amsthm}
\usepackage{booktabs}
\usepackage{graphicx}
\usepackage{xcolor}
\usepackage[hidelinks]{hyperref}
\usepackage[numbers,sort&compress]{natbib}

\newcommand{\tbd}[1]{\textcolor{red}{[TBD: #1]}}
\newcommand{\sys}{\textsc{wardowl}}
\newtheorem{definition}{Definition}
\newtheorem{observation}{Observation}

\title{Leak-Free but Not Loss-Free:\\Verifying Permission-Aware Retrieval for RAG\\across Vector Stores and Authorization Engines}
\author{Anonymous Authors}
\date{}

\begin{document}
\maketitle

\begin{abstract}
Retrieval-augmented generation (RAG) places the chunks most similar to a question into a language
model's prompt. Similarity ignores authorization, so any index shared by readers with different
permissions can hand a user's question a chunk that user may not read; OWASP lists this as LLM08.
The remedy is old: push the access predicate into the query and check what comes back. What is not
known is whether today's stack implements it faithfully. We study permission-aware retrieval as
deployed: an enforcement layer that compiles an organization's existing authorization policy into
filters for production vector stores, re-checks every hit against the authorizer, and fails closed.
We contribute (i) a formulation of secure, fresh, sound, complete and count-honest top-$k$ retrieval
under an external authorizer, with a stated trusted computing base; (ii) differential conformance
testing of filter semantics across 13 of the 14 vector stores the layer supports and seven
authorization engines, which found two undocumented store behaviors that cost recall, two adapter
faults and documented limits that fail closed, and confirmed a type gap in the layer's Cerbos
translator; (iii) a measurement of disagreement between authorization engines' filter compilers and
their own check APIs, which at one revision was confined to that type gap, and of revocation delay per
path, up to the cache's time to live or the engine's staleness; (iv) an audit of retrieval quality
under permission-shaped filters, where mean recall@10 at common settings ranged from 0.55 to 0.98 over
eight stores searched approximately, and one store returned no results at all to some least-privileged
people in five of eight runs, which the evaluated release reported as complete; (v) detection and
bounded repair of such under-retrieval; (vi) a measurement of side channels in filtered retrieval; and
(vii) an end-to-end evaluation with canary documents over 107,540 chunks, 20 tenants and 1,105
questions per run, including red-team prompts: four runs of each of eleven configurations, with no
leak and a one-sided 95\% bound of $6.8 \times 10^{-4}$ leaks per question for each configuration and
$1.6 \times 10^{-2}$ per person, over the generator's question distribution. No returned chunk was
unreadable at the revision the re-check read, in any tested configuration; information about
unreadable chunks still was revealed through latency, counts, recall and scores, and retrieval
quality varied widely.
\end{abstract}

\section{Introduction}\label{sec:intro}
% Introduction: body of Section~\ref{sec:intro} (the \section heading and its label stay in
% main.tex). Every number is from the section it summarizes.

A retrieval-augmented generation (RAG) assistant answers from the chunks of an index that are most
similar to the question~\cite{lewis2020retrievalaugmented}. In an organization those chunks come
from documents with different readers: a payroll file, a board deck, one customer's contract and a
public handbook can share one index. Similarity knows nothing of who may read what, so without
enforcement the chunk most similar to an employee's question may be one that employee may not open,
and the model will repeat it. OWASP lists this as LLM08, vector and embedding
weaknesses~\cite{owasp2025llm082025}, and revoked documents have stayed retrievable through a cached
index~\cite{roychowdhury2024confusedpilot}.

The remedy is old. Query modification adds the access predicate to the query before it
runs~\cite{stonebraker1974access}; multi-user full-text search showed that statistics computed over
unreadable documents reveal them~\cite{buttcher2005security}, and later attacks turned that observation, and
timing, into term and plaintext recovery against shared search indexes and row-level
security~\cite{wang2017sidechannel,espiritu2026plaintext}; enterprise search checked results late, in
growing blocks~\cite{bailey2006secure}; a search over content that Zanzibar protects checks its candidate results
against a central authorization service~\cite{pang2019zanzibar}. Managed RAG services combine these
steps~\cite{aws2026access}, and recent work builds permissions into the vector index
itself~\cite{zhong2026honeybee,han2026dont}.

Three things stand between an organization's policy and the chunks a model sees. \emph{First, the policy is
compiled by translators.} Filter APIs are code paths separate from check APIs, every store has its own filter
language, and we found no published guarantee that the paths agree; they have diverged, as when SpiceDB's
LookupResources missed results its other APIs computed correctly~\cite{nvd2025cve202565111}. \emph{Second, the
filter runs inside an approximate index.} Recall in filtered nearest-neighbor search depends on a filter's
selectivity and its correlation with the question, though neither alone predicts
it~\cite{patel2024acorn,lim2026revisiting}; a permission filter can be selective, takes a disjunction over many
groups, and admits chunks by organizational structure, not meaning. \emph{Third, failures are
silent.} Under Truman semantics~\cite{rizvi2004extending} a subject who receives three chunks, or none,
cannot tell whether that is all it may read or all that the index found.

This paper is a measurement study of one enforcement layer, \sys{}, an open-source library that the author develops and
whose public repository was created in October 2026. It compiles an organization's policy into filters for 14 vector
stores, from ACL metadata or from four authorization engines, checks hits through those and three more engines, re-checks
every hit and fails closed (Section~\ref{sec:system}). The results describe this layer at the versions we state, on
synthetic permission layouts. We ran no adaptive adversary and compare no other layer; our one native baseline is a side
run on another corpus (Section~\ref{sec:qual-native}). Under stated assumptions, the
re-check rests membership in the readable set on a trusted computing base around the check API
(Observation~\ref{prop:reduction}), whose size we did not measure: a faulty
filter or a lossy index can then cost recall but cannot put a chunk outside the readable set into an answer. We ask how much is lost, where, how fresh the check is, what the answers still reveal, and whether anyone can tell.

We found the following. Outside positive controls built to leak, in what we tested no returned chunk lay outside the subject's readable set at the revision the
re-check read; we call such a chunk a \emph{leak}. In the end-to-end runs the re-check dropped nothing, so those runs would
look the same with it off; only a one-seed mutation control, and a data-shape profile of one translator, exercised it. The
revision can be stale: revoked chunks were served for up to 38.4\,s after an invalidation, past the stated
bound. Other observables of an answer revealed information about unreadable chunks. And on our
layouts retrieval quality ranged from near perfect to missing most of the exact top ten for some people, unseen by the application.

\paragraph{Contributions.}
\begin{itemize}
  \item \emph{A formulation} (Section~\ref{sec:problem}). We adapt the sound and maximal criteria of Wang et
        al.~\cite{wang2007correctness} to top-$k$ retrieval under an external authorizer, replace their indistinguishability
        notion of security with a weaker property of the returned set, add freshness, with a measured staleness $\delta$
        per path, and count-honesty for the claim that an answer is complete, and state the trusted computing base.
  \item \emph{Testing of the enforcement path} (Sections~\ref{sec:conformance},
        \ref{sec:engines} and~\ref{sec:e2e}). Differential testing across 13 of the 14 stores found no hit outside the
        reference on its generated data and no altered metadata; two truncations, which at long values lose matching chunks and, as separate reproducers show, also admit a longer value by its prefix; two availability faults
        in the layer's adapters; and limits that fail closed, not all of them documented. Only four of the seven engines have filters to compare.
        At one revision, over 6.3 million user-document pairs from one generated model per
        profile, filters and checks never disagreed in the profiles that hold the model as intended. No chunk outside the readable set was returned in 1,050 fault episodes (a further 40 control episodes turned denials into allows and leaked, as they must) or
        in four runs of each of eleven end-to-end configurations ($8.0 \times 10^{-4}$ per question, $1.6 \times 10^{-2}$ per person and
        0.21 per role type, one-sided 95\%, under independence assumptions Section~\ref{sec:e2e-safety} qualifies). A mutation control on one seed and an exact store shows that the leak checks fail when
        the filter is broken on purpose and that the re-check contains that fault.
  \item \emph{An audit, diagnosis and repair of retrieval quality} (Sections~\ref{sec:quality}
        and~\ref{sec:repair}). Mean recall@10 ranged from 0.55 to 0.98 over eight stores searched approximately, on
        one synthetic corpus with 19 near-identical copies of each company document, under this layer's schema and at settings some of them tuned.
        A diagnosis grid with another question mix gave Qdrant 0.656 in the shared layout and 0.952 with per-tenant graphs, consistent with
        an islanded permitted subgraph that we did not measure directly; plans recorded under an exploratory hypothesis tie
        pgvector's two modes to PostgreSQL's plan cache. A count flagged all 2,657 truly short answers
        without a false alarm, which on unchanging data follows from what a count is, but no count sees a full answer of the wrong chunks.
  \item \emph{Side channels of filtered retrieval} (Section~\ref{sec:side}). In a pre-registered forced-choice game, shown one
        query's observation in two worlds and told that in one the query's $m$ nearest non-hub passages are unreadable chunks (the worlds differ in 1,233 such chunks, the others beyond the post-filter's cap of 1,000 candidates for that query), a subject picked
        that world from the widening post-filter's latency with paired AUC 0.66 to 0.98 at $m = 100$ and 1.000 at $m = 1{,}000$ when
        it could read 10\% or 50\% of the background; reading 1\%, it had no forward latency channel (0.03 to 0.59) but a
        count channel at $m = 1{,}000$ (1.000). BM25 and IDF scores gave 0.98 to 1.00 at those $m$, in worlds that differ by 1,233 chunks whatever $m$, a known channel that prior work takes to term and text
        recovery~\cite{wang2017sidechannel,espiritu2026plaintext}. Of
        the game's hypotheses H1 and H2 were falsified, and H3b, H4, H5a and H6 are inconclusive as printed.
\end{itemize}

We do not propose an index, and we do not defend against prompt injection in readable content. The library,
the conformance kit, the pre-registrations and raw outputs, with the gaps that Section~\ref{sec:limits} lists, are released at
\url{https://github.com/muzakkirhussain011/wardowl-paper-artifact}. Appendix~\ref{app:map} maps every experiment to its versions and
runs, and Appendix~\ref{app:ledger} lists every hypothesis with its registration and verdict; labels such as P5 below refer to that ledger.


\section{Background and Problem Statement}\label{sec:problem}
% Background and Problem Statement: body of Section~\ref{sec:problem} (the \section heading and its
% label stay in main.tex). Labels defined here: sec:prob-setup, sec:prob-criteria, def:criteria,
% prop:reduction (an observation since the 2026-10-30 revision; the label is kept because other
% sections cite it), sec:prob-threat, sec:prob-metrics.

\subsection{Permission-aware top-$k$ retrieval}\label{sec:prob-setup}

\paragraph{Retrieval in RAG.}
A retrieval-augmented generation (RAG) pipeline embeds a question, retrieves the $k$ stored chunks
most similar to it, and places them in a language model's prompt~\cite{lewis2020retrievalaugmented}.
Let $D$ be a corpus of chunks, each with an embedding $\mathbf{x}_d$ and metadata $m_d$, and let
$\sigma(q,d)$ be the similarity (cosine or inner product) between a query embedding $q$ and
$\mathbf{x}_d$. For $X \subseteq D$, $\mathrm{top}_k(q,X)$ denotes the $\min(k,|X|)$ chunks of $X$ most
similar to $q$, ties broken by a fixed order on chunk identifiers. Vector stores answer this query
approximately, with indexes such as HNSW~\cite{malkov2020efficient}, DiskANN~\cite{subramanya2019diskann}
or IVF~\cite{douze2024faiss}; a filtered query $\mathrm{search}(q,\varphi,k)$ restricts the answer to
the chunks $E(\varphi) = \{d \in D : m_d \models \varphi\}$ whose metadata satisfy a predicate
$\varphi$, filtering before, after or during index traversal~\cite{patel2024acorn,chronis2025filtered}.
Similarity ignores authorization: when readers with different permissions share one index, the chunk
most similar to a question may be one the asker may not read~\cite{arceo2026securing,owasp2025llm082025}.

\paragraph{Subjects, the policy and the authorizer.}
A \emph{subject} $s$ is the principal on whose behalf a query runs; an agent acting for a user is
authorized as that user, never with permissions the user lacks~\cite{jones2020oauth,mitre2026atlas}.
Permissions are stated by a \emph{policy}. Given the authorization data at revision $\tau$ (a
snapshot of relationships, ACLs and policy text, committed at time $t(\tau)$), the policy decides
$P_\tau(s,d) \in \{0,1\}$, whether $s$ may read chunk $d$. The \emph{readable set} of $s$ is
$R_\tau(s) = \{d \in D : P_\tau(s,d) = 1\}$, and the \emph{exact answer} to $q$ is
$T^\star = \mathrm{top}_k(q, R_\tau(s))$. Our experiments compute $P$ with an independent oracle
written without the layer's code (Section~\ref{sec:prob-metrics}). The policy is enforced by an
external \emph{authorizer}: a relationship-based engine in the style of
Zanzibar~\cite{pang2019zanzibar}, a policy engine, or ACL metadata copied from a source system. Its
\emph{check API} answers one question at a time, $\mathrm{check}_\tau(s,d) \in \{0,1\}$, reading
revision $\tau$~\cite{pang2019zanzibar,authzed2026consistency}; it is the implementation of $P$ that
an enforcement layer trusts. Many authorizers also expose a \emph{filter API}, such as OpenFGA's
ListObjects, SpiceDB's LookupResources, OPA's partial evaluation or Cerbos's query
plans~\cite{openfga2026perform,authzed2026protecting,opa2026evaluating,cerbos2026cerbos}, whose output
can be translated into a store predicate $\varphi_s$. The filter API and the check API are separate
code paths, and we found no published guarantee that they agree (Section~\ref{sec:related}).

\paragraph{The enforcement layer.}
An enforcement layer sits between the RAG application and the store. For each request it
(i)~takes $s$ from a verified token; (ii)~obtains $\varphi_s$ when a filter can be compiled;
(iii)~searches the store with $\varphi_s$ for at least $k$ candidates, or without a filter when none
can be compiled; (iv)~re-checks every candidate with $\mathrm{check}_\tau$, dropping those that fail
and counting them as \emph{violations}; (v)~widens the search until $k$ candidates survive or a cap is
reached; and (vi)~fails closed, returning nothing, when the authorizer cannot answer. It returns a
ranked list $H$ of at most $k$ chunks and a flag $c \in \{0,1\}$ claiming that $H$ is complete.
Section~\ref{sec:system} describes the layer we test, \sys{}, and credits the precedents of each step;
here we fix what such a layer should guarantee.

\subsection{Correctness criteria}\label{sec:prob-criteria}

Wang et al.\ judge fine-grained access control in relational databases by three criteria: an answer
is \emph{sound} if it is consistent with the answer computed without access control, \emph{secure}
if it reveals nothing the policy forbids, and \emph{maximal} if it returns as much as the first two
allow~\cite{wang2007correctness}. A recent vision paper asks for these criteria to be reinterpreted
for vector search~\cite{yalamarthi2026policyaware}. We adapt them to top-$k$ retrieval under Truman
semantics, in which a query is answered over the subject's restricted view rather than
rejected~\cite{rizvi2004extending}. Under Truman semantics the reference answer is the exact answer
over the readable set, so soundness becomes a property of the returned chunks relative to $T^\star$,
and maximality splits into two failures that need different detectors.

\begin{definition}[Criteria for permission-aware top-$k$ retrieval]\label{def:criteria}
Fix a subject $s$, a query $q$ and a depth $k$, for a request issued at time $t$. Each hit $d$ is
checked at some revision $\tau_d$ (different hits may be checked at different revisions, for example
when some decisions come from a cache), and $\tau$ is the latest of these revisions. $\varphi_s$ is
the predicate the store searched under (the constant true for an unfiltered search), $\sigma_{(k)}$
the $k$-th highest similarity in $R_\tau(s)$ (the lowest if $|R_\tau(s)| < k$), and
$T^\star_\epsilon = \{d \in R_\tau(s) : \sigma(q,d) \ge \sigma_{(k)} - \epsilon\}$ the chunks that
belong to the exact answer up to ties within $\epsilon$. A result $(H,c)$ is
\begin{itemize}
  \item \emph{secure} if $d \in R_{\tau_d}(s)$ for every $d \in H$: no returned chunk lies outside
        the readable set at the revision its check read;
  \item \emph{$\delta$-fresh} if it is secure and every $\tau_d$ includes every write committed before
        $t - \delta$, so that a permission revoked more than $\delta$ before the request, and not
        granted again, admits nothing in it;
  \item \emph{sound} if $H \subseteq T^\star_\epsilon$ (mis-retrieval, returning chunks that are not
        among the exact answers, violates it);
  \item \emph{complete} if $|H| = \min(k, |R_\tau(s)|)$ (under-retrieval violates it);
  \item \emph{maximal} if it is sound and complete, that is, $H$ is $\mathrm{top}_k(q, R_\tau(s))$
        up to ties within $\epsilon$; and
  \item \emph{count-honest} if $c = 1$ implies $|H| = \min\bigl(k, |R_\tau(s) \cap E(\varphi_s)|\bigr)$.
\end{itemize}
\end{definition}

Security, and freshness for the $\delta$ that the serving path states, are required of every result,
and count-honesty is required of the flag. Soundness and completeness are what approximate search
trades away; permission-recall@$k$, the under-retrieval rate and the false-completeness rate
(Section~\ref{sec:prob-metrics}) measure the shortfall. Count-honesty is relative to $\varphi_s$: it
implies completeness only when $\varphi_s$ has no false denials, $R_\tau(s) \subseteq E(\varphi_s)$,
and it never implies soundness, since a full answer of the wrong chunks satisfies it. A flag with
$c = 1$ only for maximal results would need exact search to set; Section~\ref{sec:sys-complete} shows
that the layer we test is count-honest relative to $\varphi_s$ and no more. Wang et al.'s soundness
compares an answer with the unrestricted one; ours compares it with the exact answer over the
readable set, which is what a Truman subject is owed. Score fidelity is a separate conformance
property (finite scores in non-increasing order, Section~\ref{sec:conformance}); mis-ordering inside
a store, such as the slightly out-of-order output of pgvector's relaxed iterative
scans~\cite{pgvector2026pgvector}, shows in recall, not in security. Security concerns the returned
set only. B\"{u}ttcher and Clarke's stricter requirement,
that no aspect of the output depend on unsearchable content~\cite{buttcher2005security}, is a
noninterference condition~\cite{goguen1982security} on everything a caller observes, including timing
and the flag; we treat those observables as side channels (Section~\ref{sec:prob-threat}). Under
Truman semantics a short answer looks like an answer to a subject who can read little; the flag $c$
is a partial step toward non-Truman behavior, letting the layer say that an answer may be
incomplete.

\paragraph{Security does not depend on the filter or the store.}
We make three assumptions.
\begin{itemize}
  \item[(A1)] \emph{Correct checks}: $\mathrm{check}_\tau(s,d) = P_\tau(s,d)$ for every subject, chunk
        and revision $\tau$ the check reads.
  \item[(A2)] \emph{Faithful metadata}: for every candidate the store returns, the identifier and
        attributes that the check evaluates belong to the chunk whose content is returned.
  \item[(A3)] \emph{Fail closed}: an error, timeout or malformed answer from the authorizer ends the
        request with no hits.
\end{itemize}

\begin{observation}[Security does not depend on the filter or the store]\label{prop:reduction}
Under \textup{(A1)} to \textup{(A3)}, if the layer returns only candidates that passed
$\mathrm{check}$, every result it returns is secure, whatever predicate $\varphi_s$ was compiled and
however the store evaluated or approximated it. If, moreover, every check reads a revision that
includes every write committed before $t - \delta$, every result is $\delta$-fresh.
\end{observation}

\begin{proof}[Argument]
Each $d \in H$ passed $\mathrm{check}_{\tau_d}(s,d) = 1$ on the identifier and attributes the store
returned with it. By (A2) these describe the returned content, and by (A1) $P_{\tau_d}(s,d) = 1$, so
$d \in R_{\tau_d}(s)$. If a check cannot be completed, (A3) gives $H = \emptyset$. The freshness part
restates the condition on $\tau_d$. No step refers to $\varphi_s$ or to the store's search procedure.
\end{proof}

With the readable set defined by the policy and (A1) equating the check with the policy, the
statement is close to definitional, and we call it an observation for that reason. Its content is
what it leaves out: no assumption about $\varphi_s$, the filter APIs, the translators, the store's
filter semantics or its index. It moves the burden of security onto a small trusted computing base,
stated below, and turns everything else into a question of completeness, soundness and freshness. It
also holds for every query, chosen adaptively or not, which no sampled measurement can show.

Three consequences follow. \emph{First, the re-check sees filter faults in one direction only.} A
\emph{false admit}, where $\varphi_s$ accepts a chunk outside $R_\tau(s)$, reaches the re-check, which
drops it and counts a violation; the violation count thus measures online how far the pushed-down
filter over-approximates. A \emph{false denial}, where $\varphi_s$ rejects a readable chunk, never
becomes a candidate and silently costs completeness. Truncated allow-lists (OpenFGA's non-streaming
ListObjects stops at 1,000 results or 3 seconds by default~\cite{openfga2026openfga}), divergent store
filter semantics and approximate-index misses all fail this way, and only an offline oracle finds
them: differential conformance testing against a reference semantics
(Section~\ref{sec:conformance}) or an audit against exact search over $R_\tau(s)$
(Section~\ref{sec:quality}). \emph{Second, the assumptions mark where security can still fail.} (A1)
excludes wrong upstream permissions and, when the check reads ACLs copied into chunk metadata, holds
only as far as the copy is fresh; (A2) fails under any store or ingestion fault that binds content to
the wrong identifier or metadata, which no re-check can detect; and security holds at the revision
the check reads, so how stale that revision may be is a separate guarantee, freshness, which we
state per path below. \emph{Third, a violation is a translation fault only at a single revision.} If
$\varphi_s$ is compiled at revision $\tau_f$ and hits are checked at $\tau_c \neq \tau_f$, a
disagreement may be a legitimate ACL change between the calls. Zanzibar's snapshot tokens and
SpiceDB's revision tokens prevent such mixing~\cite{pang2019zanzibar,authzed2026protecting}, and we pin
filter and check to one revision when measuring disagreement (Section~\ref{sec:engines}).

\paragraph{Trusted computing base.}
The observation trusts (i)~the authorizer's check API and the authorization data it reads (A1),
including the freshness of ACLs copied into chunk metadata; (ii)~the layer's mapping of a request to
checks: the subject taken from a verified token (Section~\ref{sec:sys-identity}) and the adapter code
that turns the subject, its groups and each chunk's identifier and attributes into the engine's check
request; (iii)~the binding of content to identifier and metadata in the store and in ingestion (A2);
(iv)~the decision cache, its keys, time to live and invalidation, and the engine's own staleness,
which together set $\delta$; and (v)~the fail-closed path (A3). It does not trust the filter APIs, the
translators, the store's filter semantics or index, the widening loop, or a reranker or redactor,
which receive only permitted hits. The trusted part has failed before. Two faults in the layer's own
adapters let the check itself admit, beyond any re-check: a SpiceDB adapter that checked groups as
separate subjects and so bypassed an exclusion, and a Cedar adapter that dropped values a
\texttt{forbid} rule needed (Table~\ref{tab:conf-faults}); the project's acceptance scenarios and
review found them, and both were fixed before the version we evaluate. With the re-check turned off,
the filter and its translator join the trusted computing base (Section~\ref{sec:sys-retrieval}).

\paragraph{Caches and staleness.}
Let the engine answer from a snapshot that may lack the writes of the last $\delta_e$ seconds, and let
the layer cache each decision for at most a time to live $\ell$. A decision used at time $t$ was then
read at a revision that includes every write committed before $t - (\ell + \delta_e)$, so under (A1)
to (A3) the layer is $(\ell + \delta_e)$-fresh, or $\delta_e$-fresh without a cache. An explicit
invalidation is meant to restore $\delta_e$-freshness for the subjects it names; it can do so only if
the cache does not then store answers the engine computes from a snapshot older than the
invalidation. Section~\ref{sec:eng-revoke} measures $\delta$ per path (Table~\ref{tab:eng-revoke}).
When permissions are copied into chunk metadata, and for engines read at full consistency, revocation
took effect in a median of 0.01\,s or less; for engines behind a cache it took the time to live
(medians within 34\,ms of 5\,s and 30\,s), and invalidation ended that delay at once, or within 1\,s
on a replica sharing the cache. Two sources of truth were stale themselves. Under SpiceDB's default
\texttt{minimizeLatency}, revoked chunks were served for a median 3.3\,s and, in one of 40 trials of
the first run, 8.7\,s, beyond the 6.5\,s we had registered as the limit (the documented 5.5\,s plus
1\,s for polling); a rerun stayed within 5.2\,s. Cerbos took a median 2.5\,s to reload a rewritten policy file, and
states no bound. In front of these two engines, the release we evaluated violated its stated bound:
after an invalidation the next call refilled the cache from the stale snapshot and held it for the
full time to live, so 79 of 80 SpiceDB revocations with invalidation exceeded the bound, with medians
of 30 and 31\,s. The fix caches nothing for the engine's declared staleness after an invalidation, and
a rerun confirmed it: all 160 SpiceDB and Cerbos revocations with invalidation took effect within the
declared staleness plus 1\,s, and 1\,s more on the replica.

\paragraph{Count certificates.}
Security needs no knowledge of $R_\tau(s)$; completeness does. A layer can set $c = 1$ count-honestly
only if (a)~the store searched exactly under $\varphi_s$, deep enough that $k$ candidates survived the
re-check or none remained, or (b)~a filter-only count shows that every chunk passing $\varphi_s$ was
among the candidates. Conversely, a count of at least $k$ alongside fewer than $k$ candidates proves
under-retrieval; OpenSearch uses this test to switch to exact search~\cite{opensearch2026efficient}.
Both certificates are relative to $\varphi_s$, so completeness with respect to $R_\tau(s)$ also needs
a filter free of false denials, which we test offline rather than certify at run time. A layer that
marks a short or empty approximate result complete is not count-honest; the release of \sys{} we ran
end to end was not, and in the runs of Section~\ref{sec:e2e} it reported empty answers as complete for
up to ten least-privileged people per run, each of whom could read 624 chunks.

\subsection{Threat model}\label{sec:prob-threat}

\paragraph{Adversary.}
The adversary is an authenticated subject, or an agent or application acting for one, who may issue
any number of arbitrary queries through the layer, chosen adaptively and including adversarial and
red-team prompts, to obtain the content of any chunk outside the subject's readable set, across
users or tenants, or to learn whether such chunks exist near a query. The adversary does not control
the authorizer, the store or the layer's host. Observation~\ref{prop:reduction} covers adaptive
queries; our measurements sample questions from a generator and bound failure rates over that
distribution only (Section~\ref{sec:prob-metrics}).

\paragraph{In scope.}
(i)~Cross-user and cross-tenant disclosure through retrieval results. (ii)~Faults in the composed
components: errors in an authorizer's filter compilation; filter-semantics bugs and gaps in stores,
such as type coercion, Unicode and case handling, and nested field names; and approximate-index
behavior under selective or misaligned filters. (iii)~Authorizer outages, slow answers and malformed
answers, which must fail closed. (iv)~Identity confusion: a retrieval tool that takes a user
identifier as an argument lets a prompt-injectable model choose whose permissions apply, a confused
deputy~\cite{hardy1988confused}; the layer takes the subject only from a verified token, following
RFC 8693 and the MCP authorization specification~\cite{jones2020oauth,mcp2026model}.
(v)~Revocation: a permission removed after indexing must stop admitting chunks within the $\delta$ the
serving path states. Context added after the vector stage, such as parent documents, neighboring
chunks or graph hops, is covered only if it passes through the same re-check~\cite{thornton2026retrieval}.

\paragraph{Out of scope.}
(i)~Prompt injection carried in readable content, poisoning of readable content~\cite{zou2025poisonedrag},
and anything the model does with retrieved content, including exfiltration after correct retrieval
as in EchoLeak~\cite{nvd2025cve202532711}; we report that the model obeyed every injected notice it
retrieved (Section~\ref{sec:e2e}) to mark where our claim ends. (ii)~Wrong upstream permissions:
over-broad ACLs are enforced faithfully, as by assistants that inherit source
permissions~\cite{microsoft2026data}. (iii)~Extraction or membership inference over readable
content~\cite{zeng2024good,naseh2025riddle}. (iv)~Infrastructure: compromise of the store, its
credentials, backups or hosts. Because embeddings can be inverted~\cite{morris2023text,zhang2025universal},
our guarantee covers the query path only; the index needs storage-level protection, and the layer
never returns or logs raw vectors, query embeddings included. (v)~Several participants per
interaction, where content must be authorized for each~\cite{bhatt2025enterprise}; we model one
subject per request.

\paragraph{Side channels.}
Observation~\ref{prop:reduction} says nothing about what the timing or shape of an answer reveals.
Widening rounds, latency and the flag $c$ depend on how many unreadable chunks lie near the query, and
a filtered graph traversal may depend on the unreadable nodes it passes. Hybrid scoring that mixes
BM25 or IDF-weighted sparse scores, or normalizes or fuses scores over candidates that include
unreadable chunks, reopens B\"{u}ttcher and Clarke's collection-statistics
channel~\cite{buttcher2005security}; the dense similarity of one query and one chunk does not, because
it ignores collection statistics. We do not assume these channels away: Section~\ref{sec:side}
measures them, and we say that such a channel \emph{reveals} information about unreadable chunks,
reserving \emph{leak} for returned content.

\subsection{Metrics}\label{sec:prob-metrics}

Metrics use the policy oracle, which evaluates $P$ without calling the layer, and exact search over
$R_\tau(s)$ by exhaustive scoring. \emph{Selectivity} is the fraction of the corpus that passes a
filter ($|R_\tau(s)|/|D|$ for the permission filter), following ACORN~\cite{patel2024acorn}; some
systems use the complement~\cite{wang2021milvus}.

\paragraph{Leaks, and bounds per question and per subject.}
A \emph{leak} is a returned chunk outside $R_\tau(s)$ as the oracle decides it, or a canary string from
an unreadable document found in a hit, a prompt or an answer; canaries follow the Secret Sharer's
inserted secrets~\cite{carlini2019secret}, used as a deterministic enforcement check rather than a
memorization estimate. When no leak occurs, the unit of the bound matters. Hits are not independent
trials: they cluster by question and by subject, and a fault in the layer would affect a question or
a permission shape, not a random hit. We therefore bound the leak rate per question and per subject,
with the one-sided 95\% bound $1 - 0.05^{1/n} \approx 3/n$ for zero events in $n$ trials, the rule of
three~\cite{hanley1983if}. For the end-to-end runs of Section~\ref{sec:e2e} this gives about
$6.8 \times 10^{-4}$ per question for each configuration (4,420 questions) and $1.6 \times 10^{-2}$
per subject (the scenario's 189 people, who recur across runs), not the $3.8 \times 10^{-6}$ that
783,888 checked hits would suggest; for the injected faults of Section~\ref{sec:eng-revoke}, about
$2.8 \times 10^{-3}$ per fault episode (1,050 episodes). These bounds hold over the distribution of the
question generator and of the injected faults; they say nothing about an adaptive adversary, for
which only Observation~\ref{prop:reduction} and its trusted computing base speak. The
\emph{violation rate}, the share of store candidates the re-check rejects, measures how far the
pushed-down filter over-approximates, not the security of the layer.

\paragraph{Permission-recall@$k$.}
For a question whose subject can read something, permission-recall@$k$ is
$|H \cap T^\star_\epsilon| / \min(k, |R_\tau(s)|)$, with $T^\star_\epsilon$ as in
Definition~\ref{def:criteria} and $\epsilon = 10^{-4}$, so that ties do not penalize a store. Measured against exact search
over the readable set rather than the store's filtered set, it charges filter false denials and index
misses alike; recall against the store's own exact search under the same filter separates the two.

\paragraph{Under-retrieval and false completeness.}
The \emph{under-retrieval rate} is the share of retrievals that are not complete, returning fewer than
$\min(k, |R_\tau(s)|)$ hits, and the \emph{empty-result rate} the share returning none to subjects who
can read something. A count can detect under-retrieval, but mis-retrieval, where $k$ hits come back
but not the right ones, shows only in recall; in the enterprise scenario it, not empty results, caused
most of the recall loss (Sections~\ref{sec:quality} and~\ref{sec:repair}). We report the confusion
matrix of $c$ against the oracle; the \emph{false-completeness rate}, the share of results with
$c = 1$ that are not maximal, measures how far a count-honest flag stays from certifying maximality.

\paragraph{Latency.}
End-to-end latency per retrieval (filter compilation, search, re-check and widening) is reported at
p50, p95 and p99, each the median of three passes after warm-up, with thread count, batch mode and
cache state stated; widening rounds and authorizer calls are recorded for the side-channel analysis.

\paragraph{A measured correlation axis.}
Recall under a filter depends on selectivity and on where the readable set lies relative to the
query. Prior work calls the second property query correlation and controls it
explicitly~\cite{patel2024acorn,zhang2026vecbench,amanbayev2026filtered,lim2026revisiting}; our
``aligned'' and ``misaligned'' permission layouts are positive and negative instances of it, not a
new axis. We therefore place each (question, subject) pair on measured axes: global selectivity;
ACORN's query correlation, which compares the query's distance to the readable set with its expected
distance to a random set of equal size~\cite{patel2024acorn}; Global-Local
Selectivity~\cite{amanbayev2026filtered}; HCBGen's $\alpha$-hardness~\cite{lim2026revisiting}; the
local selectivity $\rho_m$, the share of readable chunks among the $m$ nearest unfiltered neighbors
($m \in \{100, 1000\}$); and the unfiltered rank of the first readable chunk. Because selectivity and
correlation alone are unstable proxies for difficulty~\cite{lim2026revisiting}, we report results
against several of these. Rates carry 95\% confidence intervals from a cluster bootstrap over
subjects, and comparisons between configurations are paired on the same questions.


\section{The Enforcement Layer Under Test}\label{sec:system}
% The Enforcement Layer Under Test: body of Section~\ref{sec:system} (the \section heading and its
% label stay in main.tex). Labels defined here: tab:sys-precedents, sec:sys-contracts, eq:sys-in,
% sec:sys-retrieval, sec:sys-complete, sec:sys-compose, sec:sys-identity.
% Sources: the library's retriever.py, filters.py, authz.py, caching.py, events.py, mcp.py,
% server.py, docs/security.md and CHANGELOG.md. Completeness is described as on the
% honest-completeness change (the optional exact_search backend attribute, PR #166), which the
% paper reports; the earlier behavior is stated because Section e2e ran before the change.

\sys{}, the layer we test, is a Python library that implements steps (i) to (vi) of
Section~\ref{sec:prob-setup} between a RAG application, an authorizer and a vector store.
Table~\ref{tab:sys-precedents} separates what has precedent from what we add. The enforcement
mechanism is not ours. We add a single reference semantics for permission filters, which makes the
mechanism testable across stores, and a completeness flag that does not overstate what an
approximate index returned.

\begin{table}[tbp]
\centering
\scriptsize
\newcommand{\sysrag}{\raggedright}
\setlength{\tabcolsep}{4pt}
\begin{tabular}{@{}p{0.21\linewidth}p{0.44\linewidth}p{0.3\linewidth}@{}}
\toprule
Component & Precedent & Role in this paper \\
\midrule
\sysrag Pre-filter: the access predicate pushed into the query
 & \sysrag Query modification~\cite{stonebraker1974access}; query integration in multi-user
   full-text search~\cite{buttcher2005security}; vendor pre-filter
   integrations~\cite{cerbos2026implementing,authzed2026langchainspicedb}
 & \sysrag Mechanism under test \tabularnewline
\sysrag Widening post-filter
 & \sysrag Late binding with growing candidate blocks~\cite{bailey2006secure,mangtani2006documentlevel};
   over-fetching two to three times~\cite{openfga2026rag}
 & \sysrag Mechanism under test; its rounds are a side channel (Section~\ref{sec:side}) \tabularnewline
\sysrag Re-check of every candidate
 & \sysrag Late-bound checks per result~\cite{pang2019zanzibar}; real-time verification of
   pre-filtered documents~\cite{aws2026access}
 & \sysrag Violations as an online count of filter false admits (Proposition~\ref{prop:reduction}) \tabularnewline
\sysrag Failing closed
 & \sysrag Withholding documents on transient errors~\cite{aws2026access}; an error instead of a
   partially filtered result~\cite{microsoft2026querytime}
 & \sysrag Mechanism under test \tabularnewline
\sysrag Truman answers over the readable set
 & \sysrag Truman and non-Truman semantics~\cite{rizvi2004extending}
 & \sysrag The flag $c$ as a partial step toward non-Truman answers \tabularnewline
\sysrag Three-operator filter algebra with one reference semantics
 & \sysrag Per-store filter languages
 & \sysrag \textbf{Contribution}: the oracle for differential testing across 14 stores
   (Section~\ref{sec:conformance}) \tabularnewline
\sysrag Honest completeness
 & \sysrag Exact search when fewer than $k$ approximate results return~\cite{opensearch2026efficient}
 & \sysrag \textbf{Contribution}: an exactness declaration and a flag that never reports a short
   approximate answer as complete (Sections~\ref{sec:sys-complete} and~\ref{sec:repair}) \tabularnewline
\sysrag Composition and tenant scopes
 & \sysrag Policy-combining algorithms~\cite{oasis2013extensible}; denials that override
   grants~\cite{aws2026access}
 & \sysrag Engineering; described, not claimed \tabularnewline
\sysrag Decision cache with a revocation bound
 & \sysrag Snapshot tokens against applying stale ACLs~\cite{pang2019zanzibar}; revocation defeated
   by a cached index~\cite{roychowdhury2024confusedpilot}
 & \sysrag Engineering; described, not claimed \tabularnewline
\sysrag Subject from a verified token
 & \sysrag Token exchange~\cite{jones2020oauth}; MCP authorization~\cite{mcp2026model}
 & \sysrag Adopted, not claimed \tabularnewline
\bottomrule
\end{tabular}
\caption{Components of the enforcement layer, their precedents, and their role in this paper.}
\label{tab:sys-precedents}
\end{table}

\subsection{Contracts and filter semantics}\label{sec:sys-contracts}

For a subject $s$ and an action, an \emph{authorizer} offers $\mathrm{check}(s,H)$, which every
authorizer implements and which returns one Boolean per hit, in order; $\mathrm{filter\_for}(s)$,
which returns a filter $\varphi_s$ selecting \emph{exactly} what $s$ may read, or nothing; and,
optionally, $\mathrm{scope\_for}(s)$, which returns a \emph{necessary} condition $\psi_s$, such as
the subject's tenant, that every readable chunk satisfies. Besides ACLs copied into chunk metadata,
user functions and tenants, \sys{} supports seven engines: OPA (through a filter document its policy
returns), OpenFGA, SpiceDB and Cerbos offer
filters~\cite{openfga2026perform,authzed2026protecting,cerbos2026cerbos}, and Ory Keto, Cedar and
Casbin offer checks only. A \emph{search backend} declares whether
$\mathrm{search}(q,k,\varphi)$ applies a filter inside the search, and returns at most $k$ hits,
best first, with metadata exactly as stored, since the re-check decides on it (A2). Adapters cover
14 stores (pgvector, Qdrant, Milvus, MongoDB Atlas, Weaviate, Redis, LanceDB, FAISS, Vespa,
Cassandra, Pinecone, Elasticsearch, OpenSearch and Chroma), and an in-memory backend that evaluates
the reference semantics is their oracle. One algorithm, written as a generator of operations, runs
under a synchronous and an asynchronous driver, so both share the security logic.

Filters are built from $\mathrm{In}(f,V)$, $\mathrm{And}$ and $\mathrm{Or}$. Let $\mathcal{V}$ be the
strings and integers, Booleans excluded, with strings compared by code point, without case folding
or Unicode normalization, and never equal to integers. One function, \texttt{matches}, defines, for
metadata $m$ and $V \subseteq \mathcal{V}$,
\begin{equation}\label{eq:sys-in}
m \models \mathrm{In}(f,V) \iff
\begin{cases}
\exists x \in m(f):\ x \in \mathcal{V} \wedge x \in V & \text{if } m(f) \text{ is a list},\\
m(f) \in \mathcal{V} \wedge m(f) \in V & \text{otherwise},
\end{cases}
\end{equation}
with And and Or as conjunction and disjunction. Every store translation must agree with it, and it
is the oracle of Section~\ref{sec:conformance}. Each design rule closes a leak. (i)~Empty And and Or
are rejected, since an empty conjunction is vacuously true and would allow everything.
(ii)~Values are strings or integers only: Python equates \texttt{True}, \texttt{1} and \texttt{1.0}
while JSON-based stores compare by type, so floats, nulls, Booleans, nested lists and objects never
match, and a bare string such as \texttt{"12"} is rejected rather than read as \texttt{"1"} and
\texttt{"2"}. (iii)~A missing field never matches, so a chunk without permission metadata is readable
by nobody. (iv)~A filter that can match nothing, such as an In with no values, makes translators
raise \texttt{NothingMatches} and backends return nothing without a query, because Pinecone,
MongoDB, Milvus, LanceDB, Redis and Vespa read a missing or empty filter as ``match everything''.
Parsing rejects unknown operators and nesting deeper than 100 levels.

\subsection{Retrieval}\label{sec:sys-retrieval}

By default \sys{} pre-filters when the backend supports filters and the authorizer returns
$\varphi_s$, and post-filters otherwise, inside $\psi_s$ when one is offered; a caller who demands
pre-filtering gets an error rather than a silent fallback.

\paragraph{Pre-filter and re-check.}
One search with $\varphi_s$ returns the store's answer $H_0$, truncated to $k$ so that a store that
ignores $k$ cannot widen the checks. Every hit is checked again; hits that fail are dropped, not
replaced, and counted as \emph{violations}. Because filter and check are separate code paths, a
violation is, at a single revision, a false admit by the filter or its translation, observed online
(Proposition~\ref{prop:reduction}). With the re-check off, hits are still matched against
$\varphi_s$ locally, so a store that ignores the filter still cannot leak.

\paragraph{Widening post-filter.}
Without a filter, \sys{} fetches $k \cdot o$ candidates ($o = 4$ by default), checks those not seen
before, and multiplies the fetch by $o$ until $k$ hits are permitted, the store returns fewer
candidates than requested, or a cap is reached (1,000 by default; a larger $k$ is refused). A
redactor and a reranker, if configured, receive only permitted hits and cannot add one.

\paragraph{Failing closed.}
Any exception from the authorizer ends the request with an authorization error and no hits. The
message is fixed (``authorizer failed''), the original exception kept only as its cause, because
agent frameworks show tool errors to the model and exception text can name hosts or credentials.
Malformed answers fail likewise: a filter that is not an In, And or Or, a decision list of the wrong
length, or any decision that is not a Boolean (\texttt{1}, \texttt{"no"} and \texttt{None} are
rejected, not coerced).

\subsection{Honest completeness}\label{sec:sys-complete}

Every result carries the flag $c$ of Section~\ref{sec:prob-setup}. Earlier releases set $c = 1$ for
any pre-filtered result without violations, an empty one included, and for any post-filter whose
last search came back short. In the run of Section~\ref{sec:e2e}, \sys{} thus reported empty answers
as complete for least-privileged subjects who could read hundreds of chunks; the project's notes
attribute them to pgvector's iterative scan stopping at its tuple or memory limit, a diagnosis that
Section~\ref{sec:quality} tests. The version we evaluate lets a backend declare
\texttt{exact\_search} when it ranks every chunk that satisfies the filter, so that a short answer
proves nothing else matches; only the in-memory reference and FAISS over a flat index (its default)
do. With $e$ the declaration and $v$ the number of violations, a pre-filtered result has
$c = \bigl(|H_0| = k \vee e\bigr) \wedge (v = 0)$; a post-filtered one has $c = 1$ only if $k$ hits
were permitted, or the last search came back short before the cap and either no scope applied or $e$
holds. The flag is thus honest about under-retrieval, not mis-retrieval: a full answer is taken at
the store's word, so $c = 1$ does not certify maximality (Definition~\ref{def:criteria}), and we
measure mis-retrieval offline (Section~\ref{sec:quality}). An unfiltered post-filter also still reads
a short list as the end of the store, which fails for pgvector with iterative scans off (at most
\texttt{hnsw.ef\_search} rows, 40 by default). Count-based certification and a bounded exact
fallback, as in OpenSearch~\cite{opensearch2026efficient}, are evaluated in Section~\ref{sec:repair}.

\subsection{Composition, tenants and caching}\label{sec:sys-compose}

\paragraph{Composition.}
$\mathrm{AllOf}$ admits a chunk when every part does, asking each part only about chunks that the
earlier parts admitted; $\mathrm{AnyOf}$ asks each part only about chunks still denied. Their
filters, the conjunction or disjunction of the parts' filters, exist only when every part has one.
Their scopes are asymmetric. A scope is valid when $R(s) \subseteq \{d : m_d \models \psi_s\}$.
Since $\bigcap_i R_i \subseteq R_j$ for every $j$, AllOf may conjoin whichever filters or scopes its
parts offer; since leaving a part out of a scope for $\bigcup_i R_i$ can hide readable chunks, AnyOf
offers a scope only when every part offers a filter or a scope. Every candidate is still checked, so
a wrong scope costs speed or recall, never security. The tenant authorizer admits chunks whose tenant
field equals the subject's tenant and raises when the subject has no usable tenant, so a request
that lost its tenant returns nothing rather than everything; combined with an engine that offers no
filter, it becomes the post-filter's scope.

\paragraph{Caching.}
A caching wrapper stores, per subject and action, the filter, the scope, and each decision keyed by
the chunk's identifier and metadata. Its time to live (30\,s by default) is a stated revocation
bound, and an explicit invalidation ends it at once, also for answers that calls in flight obtained
before it. Keys encode type and length, so \texttt{1}, \texttt{"1"}, \texttt{True} and \texttt{1.0}
never share an entry, and failures are never cached. Entries shared through Redis carry their
absolute expiry, so the bound holds across replicas, and, with a signing key, an HMAC-SHA256 over key
and content, so whoever can write to Redis cannot plant a decision. The bound covers the cache alone:
engine staleness~\cite{authzed2026consistency} and the freshness of ACLs copied into chunks add to
it. The cache weakens (A1) of Proposition~\ref{prop:reduction} to a decision at most one time to live
old, and its invalidation needed three successive fixes for races that stored a pre-invalidation
answer after the invalidation.

\subsection{Identity and audit}\label{sec:sys-identity}

No interface lets a client or a model choose the subject. The MCP search tool takes only a query and
a bounded result count, derives the subject from the access token the server verified for the
request, and refuses calls without one~\cite{mcp2026model}. When an agent acts for a user, the token
is issued for the user, for example by token exchange, with the agent as the client or
actor~\cite{jones2020oauth}; the agent then holds no permission its user lacks~\cite{mitre2026atlas},
and no prompt can redirect the tool to another user's permissions, the confused-deputy
pattern~\cite{hardy1988confused}. The HTTP service for other languages verifies OIDC bearer tokens,
each with its key's own asymmetric algorithm, and refuses bodies with unknown fields, so a request
cannot name a user, group or tenant. Every call, failed and cancelled ones included, emits one event
with the subject's identifier, the returned identifiers, the counts, $c$, timings and the error's
type, but never query text, query vectors or chunk text; given a key of at least 16 bytes, it adds an
HMAC-SHA256 of the query, so queries can be correlated without being stored. The violation count and
$c$ are the online signals that Sections~\ref{sec:quality}, \ref{sec:repair} and~\ref{sec:e2e}
compare with offline oracles.


\section{Differential Conformance Testing of Filter Semantics}\label{sec:conformance}
% Differential Conformance Testing of Filter Semantics: body of Section~\ref{sec:conformance} (the
% \section heading and its label stay in main.tex). Label defined here: tab:conf-findings. The kit,
% the campaign and the full fault table are in Appendices~\ref{app:conformance}
% and~\ref{sec:conf-results}, with their labels.

Observation~\ref{prop:reduction} shields security from filter faults, but not maximality, and it
rests on stores returning the metadata they stored (A2). A filter that admits too much shows up
online as violations; one that admits too little, or a store that alters metadata, does not. We
therefore test every store translation and every authorizer offline against
Equation~\ref{eq:sys-in}: differential testing in the line of the logic-bug oracles for database
engines~\cite{rigger2020testing}, with property-based generation~\cite{maciver2019hypothesis}, applied to
metadata-filter semantics across vector stores, for which we found no such
oracle~\cite{xie2025toward,wang2025towards}.

\paragraph{The kit.}
From a seed, the kit draws small corpora with confusable values (\texttt{1} and \texttt{"1"}; \texttt{"a"}
and \texttt{"A"}; \'e as one code point or two), missing and dotted fields, nested filters and subjects in 300 extra groups.
It checks searches against $E(\varphi)$: no document outside
$E(\varphi)$, $\min(k,|E(\varphi)|)$ of them, and text and metadata identical to what was stored. Stores whose
index is approximate even on these corpora (Cassandra, and Vespa and Redis with HNSW) run in a sound-only mode. For
authorizers it requires $\mathrm{check}$ to agree with $m \models \varphi_s$ on every hit, which shows consistency, not
a correct policy. Its own suite plants 18 faulty backends and six faulty authorizers that it must reject
(Appendix~\ref{app:conformance}).

\paragraph{The campaign.}
We ran the kit with fixed versions, seeds and counts against 13 stores and 13 engine targets (OPA, OpenFGA, SpiceDB,
Cerbos, Cedar and Casbin, each with two policies, and Ory Keto), beside one control of each kind. The adapters for Keto, Cedar and
Casbin offer no filter, so there the kit exercises the check path only; OPA, OpenFGA, SpiceDB and Cerbos are the engines compared with a
reference answer. Pinecone needs an API key we did not
have, so its adapter is not covered. A run is one (target, seed, variant), of 20 seeds and nine or seven data variants,
and takes the worst outcome of its examples. The first campaign completed 4,315 runs: 3,612 passed, 84 were violations and 619 were errors.
The main result is a null: both controls passed every run, no target returned a hit outside the reference on the generated data, and none altered stored metadata.
Of the violations, 39 came from the layer's Cerbos translator, whose $\mathrm{In}$ also matches a scalar where Cerbos's check denies: a type gap in the adapter, not a fault of Cerbos, which the project's review had found and the campaign
reproduced (Appendix~\ref{sec:conf-campaign}). The other 45 were Redis (26) and MongoDB (19), which cut a long term when they index it and look up the
whole term when they search: at 4,096 bytes of tag value, which the layer's hex encoding reaches at 2,048 bytes of field and
value, and at 8,191 characters. A value equal to the kept prefix of a longer one also matches it, a false admit that the re-check should catch; the reproducers, which do not use the layer, show the admit, and the campaign's generator holds no such pairs. We did not measure whether real ACL values reach such lengths. Of the errors, 479 were OpenFGA and SpiceDB refusing whole batches with HTTP~400, and
140 were stores' limits, some documented and some not, which fail loudly and return nothing.

We changed the four adapters. The second campaign, of
4,480 runs, had 3,822 passes, 39 violations, 479 refusals and 140 errors, as the diagnosis predicted, and Redis and MongoDB
passed every long- and huge-string run. The 479 refusals are 85.5\% of OpenFGA's and SpiceDB's 560 runs, and only 81 runs reached them with every
subject, so the campaign's evidence on them is thin; Section~\ref{sec:engines} tests them on data within their limits.

\begin{table}[tbp]
\centering
\scriptsize
\newcommand{\findrag}{\raggedright}
\setlength{\tabcolsep}{3pt}
\begin{tabular}{@{}p{0.17\linewidth}p{0.4\linewidth}p{0.1\linewidth}p{0.1\linewidth}p{0.17\linewidth}@{}}
\toprule
Store or engine & Finding & Found by & Effect & In \sys{}; upstream \\
\midrule
\findrag Redis 8.10; MongoDB Atlas 8.0, 8.3
 & \findrag Long values are cut when indexed (at 4,096 bytes; at 8,191 characters) and looked up whole
 & \findrag Campaign & \findrag Recall; false admit
 & \findrag Digest of long terms (PR~\#167); reported \tabularnewline
\findrag Elasticsearch, OpenSearch, Milvus, Qdrant
 & \findrag Size and syntax limits, some documented and some not, refuse a write or query
 & \findrag Campaign & \findrag Availability
 & \findrag None needed: fail closed \tabularnewline
\findrag Chroma 1.5.9, local mode
 & \findrag A flat \texttt{\$or} of about 1,000 clauses exceeds SQLite's expression depth; one of about
   16,500 or more kills the client process
 & \findrag Campaign; filter size & \findrag Availability; crash
 & \findrag Nested clauses, then local filtering (PR~\#187); reported \tabularnewline
\findrag OpenFGA, SpiceDB (adapters)
 & \findrag A subject outside the engine's limits makes it refuse the whole request
 & \findrag Campaign & \findrag Availability
 & \findrag Refused before sending (PR~\#167) \tabularnewline
\findrag Cerbos (adapter)
 & \findrag CEL \texttt{in}, \texttt{==} and \texttt{hasIntersection} translated to In, which also matches a
   scalar
 & \findrag Review; campaign & \findrag False admit
 & \findrag Translator restricted, documented; note filed \tabularnewline
\findrag pgvector 0.8
 & \findrag Iterative scans stop at either of two limits: few or no rows under selective filters
 & \findrag End-to-end runs & \findrag Recall
 & \findrag Honest completeness; limits documented \tabularnewline
\findrag pgvector via psycopg
 & \findrag The plan cache chooses at run time between an exact generic plan and lossy custom plans
 & \findrag Diagnosis grid & \findrag Recall or speed
 & \findrag None yet; rule documented \tabularnewline
\findrag SpiceDB 1.56
 & \findrag Under \texttt{minimizeLatency} a revoked permission admitted for 8.7\,s in 1 of 40 trials, unexplained (a
   rerun's 40 trials: at most 5.2\,s), and could reappear after it had gone
 & \findrag Engine study & \findrag Stale
 & \findrag Declares 11\,s (PR~\#186) \tabularnewline
\findrag SpiceDB, Cerbos (cache)
 & \findrag After \texttt{invalidate()} the cache refilled from the engine's stale answer and held the
   revoked grant for its full time to live
 & \findrag Engine study & \findrag Stale
 & \findrag Holds back for the declared staleness (PR~\#186) \tabularnewline
\findrag SpiceDB (adapter)
 & \findrag An empty LookupResources stream was read as ``nothing readable''
 & \findrag Engine study & \findrag Recall
 & \findrag No filter: post-filter (PR~\#186) \tabularnewline
\findrag OPA (partial evaluation)
 & \findrag The residual keeps \texttt{not (x in [])} even for users with no ban, and the layer's filters
   have no negation
 & \findrag Engine study & \findrag Speed
 & \findrag None; such users are post-filtered \tabularnewline
\bottomrule
\end{tabular}
\caption{This paper's findings on stores and engines. No row has been confirmed by the maintainers of the
store or engine. Table~\ref{tab:conf-faults} in Appendix~\ref{sec:conf-results} gives all 31 rows, among
them 18 that record the project's development before this paper.}
\label{tab:conf-findings}
\end{table}

Table~\ref{tab:conf-findings} collects what this paper's experiments exposed. Counted by kind, the campaign
found two store behaviors with a correctness effect, two availability faults in the layer's own adapters
and the limits of five stores, some documented and some not, which fail closed, and it reproduced one known gap. Four of the table's eleven rows
are faults in \sys{}'s own adapters or cache that a store's or engine's behavior exposed, so part of what the study found is its own
instrument; the reproducers of the reported behaviors use only the vendors' clients.


\section{Do Authorization Engines' Filters Agree With Their Checks?}\label{sec:engines}
% Do Authorization Engines' Filters Agree With Their Checks?: body of Section~\ref{sec:engines} (the
% \section heading and its label stay in main.tex). Labels defined here: sec:eng-revoke, tab:eng-revoke.
% Method, the agreement table, skew, truncation, latency and the full text on revocation and faults
% are in Appendix~\ref{app:engines}, with their labels.
% Sources: evidence/x4-x7-run1.md (GitHub Actions run 40919304635, wardowl branch paper/x4-x7 at
% e66dffc, library at fix/engine-limits e7776f0); "The fixes, tested": evidence/x7-amendment2.md (run
% 40969846593, amendment 2, library at 3c6559f).

Observation~\ref{prop:reduction} lets a pushed-down filter be wrong without a leak, but not without
cost. We ask whether the filters \sys{} compiles from four engines admit exactly what the same engines'
checks allow, separating a translation fault, which persists at one revision, from revision skew, a write
between the two calls, and from truncated list APIs. We then measure how long a revocation keeps admitting
and whether the layer fails closed. We tested OpenFGA 1.22.0, SpiceDB 1.56.2, OPA 1.21.2 and Cerbos 0.56.0
on five seeds of a generated world of 120 users, the subjects, and 1,500 documents, against an oracle that is a separate plain
implementation of the same model. \emph{One
revision} means SpiceDB's \texttt{atExactSnapshot}; OpenFGA has no snapshot token, so its store was read quiescent;
OPA and Cerbos are stateless per request (Appendix~\ref{app:engines}).

\paragraph{Agreement at one revision.}
In the seven profiles that hold the model as intended, 6.3 million (user, document) pairs produced no
disagreement between filter and check, and every check agreed with the oracle (one-sided 95\% bound on
users with any disagreement: 0.5\% per profile; X4-H1 confirmed; Table~\ref{tab:eng-agree} in
Appendix~\ref{sec:eng-agree}). These pairs repeat one generated model per profile over five seeds, and the
OPA and Cerbos profiles omit bans (Cerbos's also folders), so this has little power to find a schema-shape fault such as
CVE-2025-65111, in which SpiceDB before 1.47.1 missed LookupResources results~\cite{nvd2025cve202565111}. The only disagreement at one revision is the Cerbos type gap: when 10\% of
documents store a one-value ACL list as a plain value, the adapter's filter admits what the check denies
in 569 pairs (0.063\% of pairs; 52.5\% of users, 95\% interval 48.3\% to 56.3\%), and the re-check turns
each into a violation (X4-H2 confirmed). Neither translator dropped a negation, but vacuously: with bans no user received a
filter (X4-H3 confirmed as registered), so negation in a filter is untested. OPA's
partial evaluation keeps \texttt{not (x in [])} even for the 73 users with no ban in scope, which falsified
our availability prediction: one ban rule makes that filter unavailable to everyone, who are then
post-filtered.

\paragraph{Revision skew and truncation.}
Unpinned under a writer committing 20 toggles a second, filter and check disagreed in every OpenFGA trial,
in 81.5\% of SpiceDB trials with \texttt{fullyConsistent} (95\% interval 76.0\% to 86.5\%) and in 13.0\% of
OPA trials (8.5\% to 18.0\%). In those conditions a write between the calls explained every disagreeing pair, an attribution that is
generous where windows hold dozens of writes (medians of 59 and 208 for OpenFGA), and none reproduced; under SpiceDB's default
\texttt{minimizeLatency} it explained 226 of 228. Pinned to one SpiceDB snapshot, 200 trials produced no disagreement (X4-H4 confirmed). An online
violation is thus evidence of a translation fault only when filter and check read one revision. ListObjects and
LookupResources with the layer's limit of 1,000 truncated without a signal, and the layer gave no filter at the limit in all 540 calls (X4-H5
confirmed), at a cost: with 1,000 or 1,001 readable documents among 100,000, recall@10 fell to 0.855
to 0.891 (X4-H6 confirmed).

\subsection{Revocation and faults}\label{sec:eng-revoke}

\begin{table}[tbp]
\centering
\small
\setlength{\tabcolsep}{4pt}
\begin{tabular}{@{}lrrrrr@{}}
\toprule
Path (revocation) & No cache & ttl 5\,s & ttl 30\,s & + invalidate & shared + inv. \\
\midrule
Metadata ACL, pgvector or Qdrant & $<$0.01 & $<$0.01 & $<$0.01 & $<$0.01 & $<$0.01 \\
OpenFGA (tuple delete) & $<$0.01 & 5.02 & 30.02 & $<$0.01 & 1.01 \\
OpenFGA, check cache, \texttt{MINIMIZE\_LATENCY} & $<$0.01 & 5.02 & 30.02 & $<$0.01 & 1.02 \\
SpiceDB, \texttt{fullyConsistent} & $<$0.01 & 5.02 & 30.02 & $<$0.01 & 1.02 \\
SpiceDB, \texttt{minimizeLatency} & 3.25\textsuperscript{$\dagger$} & 5.02 & 30.02 & 30.04\textsuperscript{$\dagger$} & 31.05\textsuperscript{$\dagger$} \\
OPA (data update) & $<$0.01 & 5.03 & 30.03 & $<$0.01 & 1.02 \\
OPA (policy update) & 0.01 & 5.02 & 30.03 & 0.01 & 1.02 \\
Cerbos (policy file reload) & 2.55 & 5.02 & 30.02 & 30.03 & 31.04 \\
\bottomrule
\end{tabular}
\caption{Median seconds from a revocation to the first of three consecutive retrievals without the
revoked chunk; 40 trials per cell, the cache warmed just before. ``Shared'' polls a second replica
sharing the cache through Redis. $\dagger$: at least one trial exceeded the stated bound, the ttl
(none after invalidation, 1\,s for the replica) plus the engine's stated staleness, by more than the
1\,s we allowed for polling; Cerbos states no staleness.}
\label{tab:eng-revoke}
\end{table}

Table~\ref{tab:eng-revoke} gives the measured $\delta$ of each serving path, and a dagger marks a violation
of the $\delta$ the path states. The cache's time to live is the realized delay (X7-H2 confirmed), and invalidation ends the delay at once, or within the
shared replica's 1\,s local copy. Permissions copied into chunk metadata took effect within milliseconds of the write to the store
(X7-H4 confirmed); the delay of copying them from a source system was not measured. Two sources of truth were stale themselves. Cerbos reloaded a rewritten policy file in a
median 2.5\,s. SpiceDB under \texttt{minimizeLatency}, the layer's default, took a median
3.3\,s; in 28 polls a revoked chunk reappeared after it had disappeared, and one of 40 trials took 8.7\,s,
beyond the 6.5\,s we had registered (X7-H5 falsified), which we have not explained; a rerun's 40 trials took at most 5.2\,s.
Invalidation did not survive such an engine: the next call refilled the cache from the stale snapshot, so 79 of 80 SpiceDB revocations
with invalidation exceeded their bound (medians 30.0 and 31.0\,s, maximum 38.4\,s; X7-H1 and X7-H3 falsified in these cells), and
Cerbos's took as long.

We fixed this and four other defects the run exposed: after an invalidation the cache now caches nothing for the
engine's declared staleness (11\,s for SpiceDB without full consistency, set from the one 8.7\,s trial). A targeted rerun, an implementation check
specified in writing after the main results and before running, confirmed every prediction (A2-H1 to A2-H8): all 160
SpiceDB and Cerbos revocations with invalidation took effect within the declared staleness plus 1\,s, and 1\,s more on
the replica, at medians of 4.8\,s (at most 5.2\,s) and 3.9\,s against 30 and 31\,s before. With a maximum of 5.2\,s the rerun shows
that the fix holds, not that 11\,s is tight, and these fixes have not been through the end-to-end runs.

We injected 14 faults into each engine, and five or eight into pgvector and Qdrant, in 1,090 episodes. Outside
the control below, no retrieval returned a chunk outside the user's readable set, in 623,279 requests over 1,050
episodes. The one-sided 95\% formula gives $2.8 \times 10^{-3}$ per episode (X7-H6 confirmed), but these episodes are not samples from a
population of faults: they replay 13 fault kinds, each episode cycling 10 fixed (user, query) pairs, and the same formula over the 13 kinds gives 0.21.
The control, 40 episodes that turned the engine's denials into allows, is the fault class that does leak: it leaked in every engine, as a wrong check must, and is
excluded from the bound. Every engine-caused
failure was the layer's authorization error with a message naming no host. Recovery took a median 0.02 to
0.18\,s after a transport fault's removal and 0.08 to 0.47\,s after a killed server's restart. Two behaviors needed work: an empty LookupResources stream read as ``no
readable documents'' (4,746 retrievals returned nothing without an error, losing recall only), and store failures
surfaced as the clients' own exceptions. The rerun confirmed both fixes: in the 10,303 requests of its 20 SpiceDB episodes
no chunk leaked, and its 120 store episodes reported none either ($2.1 \times 10^{-2}$ per episode over all 140). For a connection that is accepted but never answered, only a
deadline on the whole retrieval, which the layer does not yet offer, would bound the wait.


\section{Retrieval Quality Under Permission-Shaped Filters}\label{sec:quality}
% Retrieval Quality Under Permission-Shaped Filters: body of Section~\ref{sec:quality} (the \section
% heading and its label stay in main.tex). Labels defined here: sec:qual-audit, tab:qual-x0,
% sec:qual-qdrant, tab:qual-qdrant, sec:qual-pgvector, sec:qual-frontiers, sec:qual-size.
% Sources: evidence/kaggle-scale-run-2026-10-24.md (X0), evidence/x1-underretrieval-run1.md (X1 run
% 1), investigations/filtered-recall-gaps.md (mechanisms, with source links). X1 run 2, X2 and X3
% are pending and marked \tbd{}.

Security held in every configuration of Section~\ref{sec:e2e}; this section asks what the
permission filter cost. We first audit the enterprise scenario at common settings, then diagnose the
two largest losses with pre-registered experiments, and finally map recall against throughput and
filter size. Recall is permission-recall@$k$ against exact search over the readable set
(Section~\ref{sec:prob-metrics}).

\subsection{An audit at common settings}\label{sec:qual-audit}

\begin{table}[tbp]
\centering
\small
\begin{tabular}{@{}lrrrrrr@{}}
\toprule
 & \multicolumn{3}{c}{Four runs} & \multicolumn{2}{c}{Overlap@5, first run} & \\
\cmidrule(lr){2-4}\cmidrule(lr){5-6}
Store & Recall@5 & Recall@10 & Range@10 & judy & others & p50 (ms) \\
\midrule
FAISS (flat, exact) & 1.000 & 1.000 & 0.000 & 1.000 & 1.000 & 346--368 \\
Redis & 0.983 & 0.981 & 0.012 & 0.998 & 0.982 & 28--30 \\
Elasticsearch & 0.942 & 0.949 & 0.019 & 0.987 & 0.941 & 8--9 \\
OpenSearch & 0.942 & 0.935 & 0.005 & 0.998 & 0.941 & 7 \\
Weaviate & 0.899 & 0.902 & 0.021 & 1.000 & 0.876 & 7--8 \\
Chroma & 0.877 & 0.879 & 0.010 & 0.987 & 0.876 & 1432--1481 \\
LanceDB & 0.840 & 0.874 & 0.039 & 0.949 & 0.821 & 22--25 \\
pgvector, fast mode (2 runs) & 0.695 & 0.732 & 0.000 & 0.724 & 0.686 & 6.5--8.1 \\
pgvector, exact mode (2 runs) & 0.993 & 0.993 & 0.002 & & & 386--399 \\
Qdrant & 0.557 & 0.551 & 0.026 & 1.000 & 0.523 & 4.3--4.7 \\
\bottomrule
\end{tabular}
\caption{Retrieval quality in the enterprise scenario (Section~\ref{sec:e2e}), four runs of 1,105
questions (seeds 7, 11, 13, 17). Recall is the mean over runs; Range@10 is the largest minus the
smallest recall@10; p50 is the range of per-run medians. pgvector ran in one of two modes per run
(Section~\ref{sec:qual-pgvector}); each row averages the two runs in that mode. Overlap@5 is the
share of the exact top five that a store returned, by person type, from a re-analysis of the first
run (a fast-mode run for pgvector); ``judy'' are the least-privileged people (624 readable chunks).
pgvector with OPA follows pgvector's mode in each run. Settings: Appendix~\ref{app:settings}.}
\label{tab:qual-x0}
\end{table}

Table~\ref{tab:qual-x0} shows the spread. At settings a team would plausibly ship, mean recall@10
ranges from 0.55 for Qdrant to 0.98 for Redis, and every store's loss is concentrated where the
readable set is far from the question. Most stores repeat closely: across four runs with fresh
documents and questions, no store other than pgvector moved its recall@10 by more than 0.04. Two
failures stand out. pgvector returned \emph{nothing} to some of the least-privileged people,
although each could read 624 chunks: ten in the first run, eight on red-team prompts and two on
questions about documents they could not read, and seven in another; the layer, before the change
of Section~\ref{sec:sys-complete}, called these empty answers complete. In the other two runs the
same configuration searched exactly, at fifty times the latency, and left one person without results
in one of them. Qdrant never returned a
short answer, yet for ordinary employees it returned the right documents only about half the time
(overlap@5 0.52 in the first run), while for the least-privileged people it was perfect. On targeted
questions in that run, 59.6\% of Qdrant's top five were shared FiQA passages where exact search
returned none, and the model found the target document for only 39.9\% of targeted questions
(33.8\% to 39.9\% over the four runs).

The two patterns are opposite in kind. pgvector \emph{under-retrieves}: it returns fewer than $k$
hits, which a count can detect. Qdrant \emph{mis-retrieves}: it returns $k$ hits, the wrong ones,
which only an exact baseline reveals. We pre-registered hypotheses for both before running a
diagnosis grid on Kaggle (X1; the pre-registration and its amendments are released).

\subsection{Qdrant: an islanded permitted subgraph}\label{sec:qual-qdrant}

Qdrant plans filtered search per segment from estimated filter cardinality: a filter estimated below
a threshold (about 3,333 points for 768-dimensional vectors at the default) runs as a plain, exact
scan over the payload index, and a broader one runs as filtered HNSW traversal, which scores and
expands only neighbors that pass the filter~\cite{qdrant2026search}. Extra graph edges that keep a
filtered subgraph connected are built per payload value, and only for values held by fewer points
than a cap of a few percent of the segment~\cite{couzon2026filtered}. In the scenario every tenant
value (56\% of points) and
the employee and contractor groups (65\% and 61\%) exceed the cap. Near a question about company
documents, only one of the twenty near-identical tenant copies passes the filter, so the permitted
company region is a sparse island in the graph, while the shared passages, all permitted, form a
large connected region that the traversal enters and does not leave. The least-privileged people
escape the trap because their filter is small enough to be planned as an exact scan.

\begin{table}[tbp]
\centering
\small
\begin{tabular}{@{}llrrr@{}}
\toprule
Layout & Search setting & Recall@10 & Range over builds & p50 (ms) \\
\midrule
As run (one collection) & \texttt{hnsw\_ef} 128 (as run) & 0.619 & 0.062 & 4.5 \\
 & \texttt{hnsw\_ef} 64 & 0.571 & 0.049 & 3.9 \\
 & \texttt{hnsw\_ef} 512 & 0.709 & 0.101 & 7.7 \\
 & ACORN, max selectivity 0.4 or 1.0 & 0.761 & 0.151 & 4.8 \\
 & exact & 1.000 & n/a & 35.4 \\
Tenant field marked \texttt{is\_tenant} & as run & 0.656 & 0.031 & 4.3 \\
 & ACORN & 0.830 & 0.038 & 4.5 \\
Per-tenant graphs (\texttt{is\_tenant}, $m = 0$) & as run & 0.943 & 0.009 & 3.8 \\
Shared passages in a separate collection & as run & 0.948 & 0.004 & 8.2 \\
\bottomrule
\end{tabular}
\caption{Qdrant 1.19.2 under the enterprise scenario's permission filters (X1, 1,489 questions,
three index builds per layout, means over builds). ACORN at the default gate (0.4) and at 1.0 gave
identical results. \tbd{one-segment layout after the run-2 fix.}}
\label{tab:qual-qdrant}
\end{table}

The pre-registered predictions for this hypothesis all held (Table~\ref{tab:qual-qdrant}). Recall@10
for the least-privileged people was 1.000, against 0.273 (95\% CI 0.242 to 0.305) for company
documents among everyone else; 98.6\% of the hits that did not count were shared passages; and two
interventions that reconnect the permitted region restored recall. ACORN, which explores neighbors
of filtered-out neighbors~\cite{patel2024acorn}, lifted recall@10 from 0.619 to 0.761 at an 8\%
latency cost, and per-tenant graphs~\cite{qdrant2026multitenancy} lifted it to 0.943 while also
being the fastest configuration.
Putting the shared passages in their own collection, queried separately, reached 0.948. Raising
\texttt{hnsw\_ef} from 128 to 512 added 0.090 at 1.7 times the latency, within the band we had
pre-registered as inconclusive, and did not close the gap, as expected for a traversal trapped in
the wrong region.

One prediction failed. We expected ACORN's default gate, which engages when the estimated
selectivity is at most 0.4, to engage for some people and not others, because the planner's
independence estimate for a typical employee sits near 0.4. It engaged for everyone, and the gate
at 1.0 gave identical recall. The paper therefore claims that ACORN helps but does not suffice under
this layout, not that its gate misfires. A third hypothesis, that recall varies with the segment
layout, remains open: recall@10 varied by 0.062 between three builds of the same data, but our
attempt to force a single segment did not take effect. \tbd{run 2}

For practitioners the result is a layout rule, not a parameter: when permission groups or tenants
cover large, semantically interleaved parts of a collection, give each tenant its own graph, or keep
widely shared content apart, before reaching for search-time effort.

\subsection{pgvector: the iterative scan's budget}\label{sec:qual-pgvector}

pgvector 0.8 answers a filtered HNSW query by scanning the graph in order of distance and applying
the filter to each row the scan emits; with iterative scans enabled, it continues past
\texttt{hnsw.ef\_search} until enough rows pass the filter or a limit is
reached~\cite{pgvector2026pgvector}. The limit is \texttt{hnsw.max\_scan\_tuples} (20,000 by default)
or a memory bound of \texttt{hnsw.scan\_mem\_multiplier} times \texttt{work\_mem}. A person whose
readable set is small and far from the question can exhaust the budget before any readable row is
emitted, and the query then returns nothing, which is indistinguishable from having nothing to read.

The pre-registered test of this mechanism held. Every empty answer stopped after more than half of
the budget had been spent, at a median of 20,083 rows emitted, and every setting that raised the
tuple limit to 100,000 or more removed the empty answers entirely, while raising only the memory
bound (eight times the multiplier, or 64\,MB of \texttt{work\_mem}) at the default tuple limit did
not. \tbd{X1 run 2: absolute counts; the magnitude of the X0
loss (recall@5 0.691 in the scenario run against 0.998 in the first diagnosis run) and its
explanation from investigations/x0-x1-pgvector-gap.md; ef\_search and build effects; exact
baseline.}

\subsection{Recall against throughput}\label{sec:qual-frontiers}

The enterprise scenario has one permission layout. To separate the filter's selectivity from its
correlation with the question, X2 adds a corpus whose layouts we control: the first million
passages of MS MARCO~\cite{bajaj2016ms} from BEIR, embedded with bge-small-en-v1.5 (384 dimensions,
so that ten stores fit beside the harness on a 16\,GB runner). Every passage belongs to one team of
a hierarchy of 100 divisions, 10 departments each and 10 teams each, assigned either at random
(\emph{uncorrelated}) or by hierarchical $k$-means over the embeddings (\emph{clustered}, so groups
are topical), five seeds of each. A subject receives divisions, departments and teams until it may
read a target share of the corpus, from 0.06\% to 100\%, and its filter is either one disjunction over
its groups or a tenant clause conjoined with a disjunction over a grant twice as large, a shape in
which each clause alone passes far more than both. For clustered layouts the questions are drawn
either from those whose 100 nearest passages are at least half readable (\emph{positively
correlated}) or from those with nothing readable among their 1,000 nearest (\emph{negatively
correlated}). For every question and subject we compute, from exact scores, the global selectivity,
the local selectivity among the 100 and 1,000 nearest passages, the rank of the first readable
passage, ACORN's query correlation~\cite{patel2024acorn}, the global-local selectivity of Amanbayev
et al.~\cite{amanbayev2026filtered} and the over-fetch estimate of HCBGen~\cite{lim2026revisiting}.

Ten approximate stores run against FAISS over a flat index: pgvector, Qdrant, Elasticsearch,
OpenSearch, Weaviate, Redis, Milvus, MongoDB, LanceDB and Chroma, each at its default, at five levels
of its search-effort parameter, with its store-specific options (pgvector's raised scan limits and
strict order, Qdrant's ACORN and per-tenant graphs, Weaviate's filter strategies, Redis's hybrid
policies, Elasticsearch's quantized index types) and with its own exact search where it has one, in
three fresh builds per store. Each pgvector search is planned with its own values, so the plan-cache
effect of Section~\ref{sec:qual-pgvector} cannot arise, and the share of searches whose plan uses the
HNSW index is recorded. Throughput is measured end to end through the layer with one client thread
and compared only within a job; recall is compared across jobs. The same harness runs the enterprise
scenario, whose as-run settings must reproduce the audit's recall within 0.03 for seven of nine
stores. We pre-registered four hypotheses: that negatively correlated questions lose at least 0.05
of recall@10 against uncorrelated ones at defaults in at least six stores (H2.1); that local
selectivity predicts recall better than global selectivity (H2.2); that raising search effort
narrows but does not close that gap (H2.3); and that intermediate selectivities are harder than very
small ones, which stores answer by brute force (H2.4).
\tbd{X2 full run (GitHub Actions run 40921769714): frontiers per store and regime (figure);
verdicts on H2.1 to H2.4 with intervals; recall by measured axis; validity checks V1 to V5.}

\subsection{Filter size}\label{sec:qual-size}

A person in many directory groups turns the permission filter into a long disjunction: the filter of
Section~\ref{sec:sys-contracts} holds one value per group, and enterprise users can belong to
thousands. We vary only that length. Each of five seeds builds a corpus of 100,000 vectors of 384
dimensions, the 57,638 FiQA passages embedded as in Section~\ref{sec:side-setup} and 42,362 synthetic
vectors around 200 seeded cluster centers, and makes a fixed 10\% of it readable at every level: a
readable document holds one group of each level, and the subject at level $G$ holds exactly that
level's $G$ groups, for $G$ from 1 to 40,000. The readable set, and so the exact answer, is the same
at every level; only the filter grows. Twelve stores run at default index settings, with FAISS
over a flat index and the in-memory backend as controls, and each level is timed over 30 questions
with $k = 10$ in three passes after a warm-up, a retrieval being allowed 120\,s. We record whether
each store's translator returns and how fast, whether the store accepts the filter or refuses it,
and with which error, the latency percentiles, and recall@10 against exact search over the readable
set; for pgvector also the plan. The limits the stores document that could bind are a 65,536-term
cap per \texttt{terms} query in Elasticsearch and OpenSearch, PostgreSQL's 65,535 bind parameters,
SQLite's 32,766 host parameters and expression depth of 1,000 under Chroma, and the request sizes of
MongoDB (16\,MB) and Qdrant (32\,MB); the layer's Cassandra adapter falls back to one unfiltered
query filtered locally beyond 254 groups, as documented.

We pre-registered that every translator returns within a second at every level; that at 40,000
groups Elasticsearch, OpenSearch, Qdrant, pgvector, LanceDB, Cassandra, FAISS and the in-memory
backend accept the filter and Chroma refuses it; that pgvector's median latency grows at least
fiftyfold from one group to 40,000, because above 100 values the layer sends an array test that the
GIN index cannot serve, while the controls' grows at most twofold; and that recall does not depend
on the filter's length in any approximate store that accepts it (paired difference within
$\pm 0.02$).
\tbd{X3 full run (GitHub Actions run 40940385704): acceptance per store and level with error
classes; latency against $G$; recall differences with intervals; pgvector's plans; Chroma's
process ended with a segmentation fault (signal 11) in every seed it reached, rather than refusing
the filter with an error, to be confirmed from the logs.}


\section{Detecting and Repairing Under-Retrieval}\label{sec:repair}
% Detecting and Repairing Under-Retrieval: body of Section~\ref{sec:repair} (the \section heading and
% its label stay in main.tex). Labels defined here: sec:rep-design, alg:rep, sec:rep-eval,
% tab:rep-results. Sources: the library change (PR #168, short_results), PREREG_X2_X5.md on branch
% paper/x2-x5 (X5, committed at 828ea27), results from evidence/x2-x5-run1.md (GitHub Actions run
% 40921769714, commit 884ce18), every number as x2x5_stats.py printed it.

Section~\ref{sec:quality} found two kinds of loss: \emph{under-retrieval}, where a filtered
approximate search returns fewer than $k$ hits although more readable chunks exist, and
\emph{mis-retrieval}, where it returns $k$ hits that are not the right ones. Section~\ref{sec:sys-complete}
already stops the layer from calling a short approximate answer complete. Here we go further: decide
cheaply \emph{why} an answer is short, and repair it within a stated cost. Mis-retrieval needs an
exact baseline to see, and we return to it with a sampled monitor.

\subsection{Design}\label{sec:rep-design}

\paragraph{Detection by a bounded count.}
When a pre-filtered search on a backend that does not declare exact search returns $h < k$ hits, the
layer asks the backend for $n = \min(\#\{d : m_d \models \varphi_s\},\ L)$, a filter-only count that
stops at a limit $L$. If $n \le h$, the store returned every chunk the filter admits, and the answer is
complete with respect to $\varphi_s$ (\emph{exhausted}). If $n > h$, chunks that pass the filter were
missed (\emph{missed}). The count is the second certificate of Section~\ref{sec:prob-criteria}; it
never raises a false alarm, because a count above $h$ proves a miss, and it never certifies falsely,
because a count no larger than $h$ proves that nothing else passes $\varphi_s$. Both statements hold
for the store's state at the time of the count; a write between the search and the count can make
either one stale, as any read without a snapshot can. It costs one query per short answer, and the
limit keeps that query bounded on stores that count by scanning.

\paragraph{Repair by bounded exact search.}
After a proven miss the layer can search again, exactly, under the same filter: pgvector disables the
index scan for that query, so every row that passes the filter is ranked, and Qdrant sets its
per-request exact flag. The repair runs only when the count is at most a bound $B$
(\texttt{max\_exact\_rows}, 10,000 by default), so that a subject who can read most of the corpus
cannot trigger a scan of all of it; above the bound the first hits are kept and the result says
\emph{over\_limit}. With $L = B + 1$ one count both detects and bounds. The alternative
\emph{widen} repeats the search on the same store with more effort, for example a larger scan budget
or ACORN, and merges the two answers; it is cheaper but certifies nothing, since a wider approximate
search can still miss. OpenSearch already falls back to exact search inside the engine when an
approximate filtered search returns fewer than $k$ results~\cite{opensearch2026efficient}; the layer
does the same above the store, for stores that do not.

\paragraph{Security.}
The repair is one more search step in the layer's single plan, shared by the synchronous and
asynchronous drivers. It uses the filter already obtained from the authorizer, and every hit it
returns passes the same re-check as any other. The hits of the first search that the repair
discards are re-checked in the same call, so a store that ignored the filter still shows up as
violations. A failed count or search ends the request with an error and no hits, as any other
failure does (assumption A3). Observation~\ref{prop:reduction} therefore holds unchanged.

\paragraph{A sampled recall monitor.}
Neither the count nor the flag can see mis-retrieval. For a configurable share of pre-filtered
requests, the layer also runs the exact search and reports the overlap of the two top-$k$ lists as a
metric, without returning the exact hits. This turns the offline audit of Section~\ref{sec:quality}
into a production signal, at the cost of one exact search per sampled request.

\paragraph{Reporting.}
Each result carries the outcome of detection (\emph{exhausted}, \emph{missed} or \emph{unknown}),
of repair (\emph{exact}, \emph{widened}, \emph{over\_limit} or \emph{unsupported}) and, when sampled,
the measured recall; the same fields reach the audit events and the OpenTelemetry spans and metrics
(Section~\ref{sec:sys-identity}). The default policy remains \emph{report}, which adds no query.

\subsection{Evaluation}\label{sec:rep-eval}

We evaluate the policies on pgvector and Qdrant, the two stores whose losses
Section~\ref{sec:quality} diagnosed, on the enterprise scenario (1,468 questions with something to
read) and on the 6,200 question--subject pairs of the MS MARCO layouts of
Section~\ref{sec:qual-frontiers}, on X2's prepared data and in jobs of their own, three fresh builds
per store and corpus. The base configuration is the audit's: pgvector with relaxed-order iterative
scans, \texttt{ef\_search} 100 and the default scan limits, each search planned with its own values;
Qdrant with \texttt{hnsw\_ef} 128 and no ACORN. Over the same backend and build we run
\emph{report}; \emph{count}; \emph{exact} with $B$ of 1,000, 10,000 and unbounded; \emph{widen},
whose wider search is pgvector with a million-tuple scan limit, a sixteenfold scan memory and
\texttt{ef\_search} 400, or Qdrant with \texttt{hnsw\_ef} 512 and ACORN at gate 1.0; the 1\% recall
monitor; and, for comparison, the store-native alternatives alone: pgvector's tuple limit raised to a
million, Qdrant's ACORN at gate 1.0, and, on the enterprise corpus, Qdrant's per-tenant graphs. Each
question retrieves with $k = 5$ and $k = 10$ in a first pass and twice more with $k = 10$; every hit is
re-checked by the layer and by an independent oracle, and latency is the median of the three passes
with $k = 10$.

Seven hypotheses were pre-registered. The count detects every short first search without a false
alarm (H5.1); exact repair without a bound leaves no result empty or short and reaches recall@10 of
at least 0.999 on the queries it repairs (H5.2); no repair changes the median latency of the queries that do not
trigger it by more than 5\% either way (H5.3); no repair changes Qdrant's false-completeness rate, because its loss is mis-retrieval
(H5.4); the monitor samples 1\% of searches and its recall agrees with the harness's (H5.5);
widening recovers less than exact repair on the searches that trigger it (H5.6); and on the
enterprise corpus Qdrant's per-tenant layout reaches recall@10 of 0.90 while every library arm on
the shared layout stays at or below 0.75 (H5.7).

\begin{table*}[tbp]
\centering
\footnotesize
\begin{tabular}{@{}llrrrrrrrrrrr@{}}
\toprule
 & & & & & & & & & \multicolumn{2}{c}{All (ms)} & \multicolumn{2}{c}{Triggered (ms)} \\
\cmidrule(lr){10-11}\cmidrule(lr){12-13}
Corpus, store & Arm & R@5 & R@10 & Empty & Short & Trigger & Detect & FC & p50 & p99 & p50 & p95 \\
\midrule
Enterprise, & report & 0.639 & 0.667 & 0.089 & 0.116 & 0.129 & -- & 0.551 & 17.3 & 92.5 & -- & -- \\
pgvector & count & 0.639 & 0.667 & 0.089 & 0.116 & 0.129 & 1.000 & 0.542 & 17.7 & 114.1 & 14.7 & 25.6 \\
 & exact, $B$ 1,000 & 0.747 & 0.782 & 0.000 & 0.000 & 0.129 & 1.000 & 0.480 & 17.6 & 122.7 & 25.3 & 30.6 \\
 & exact, $B$ 10,000 & 0.747 & 0.782 & 0.000 & 0.000 & 0.129 & 1.000 & 0.480 & 17.3 & 120.9 & 25.2 & 30.3 \\
 & exact, unbounded & 0.747 & 0.782 & 0.000 & 0.000 & 0.129 & 1.000 & 0.480 & 17.2 & 127.2 & 30.1 & 40.2 \\
 & widen & 0.739 & 0.775 & 0.000 & 0.000 & 0.129 & 1.000 & 0.498 & 17.9 & 410.7 & 34.5 & 360.9 \\
 & monitor 1\% & 0.639 & 0.667 & 0.089 & 0.116 & 0.129 & -- & 0.551 & 17.6 & 96.1 & -- & -- \\
 & tuple limit $10^6$ alone & 0.664 & 0.694 & 0.072 & 0.090 & 0.103 & -- & 0.543 & 17.6 & 128.9 & -- & -- \\
\midrule
Enterprise, & report & 0.658 & 0.654 & 0.000 & 0.000 & 0.014 & -- & 0.581 & 4.72 & 6.15 & -- & -- \\
Qdrant & count, exact, widen & 0.658 & 0.654 & 0.000 & 0.000 & 0.014 & n/a & 0.573 & 4.59--4.60 & 5.86--5.91 & 1.57--1.58 & 1.69--1.84 \\
 & monitor 1\% & 0.658 & 0.654 & 0.000 & 0.000 & 0.014 & -- & 0.581 & 4.71 & 6.00 & -- & -- \\
 & ACORN 1.0 alone & 0.857 & 0.855 & 0.000 & 0.000 & 0.014 & -- & 0.388 & 4.84 & 6.35 & -- & -- \\
 & per-tenant graphs & 0.938 & 0.939 & 0.000 & 0.000 & 0.014 & -- & 0.235 & 4.69 & 8.27 & -- & -- \\
\midrule
MS MARCO, & report & 0.896 & 0.900 & 0.003 & 0.017 & 0.017 & -- & 0.301 & 9.87 & 425.4 & -- & -- \\
pgvector & count & 0.896 & 0.900 & 0.003 & 0.017 & 0.017 & 1.000 & 0.301 & 9.83 & 427.7 & 21.7 & 87.7 \\
 & exact, $B$ 1,000 & 0.900 & 0.904 & 0.001 & 0.012 & 0.017 & 1.000 & 0.299 & 9.82 & 426.1 & 60.0 & 129.2 \\
 & exact, $B$ 10,000 & 0.909 & 0.915 & 0.000 & 0.001 & 0.017 & 1.000 & 0.296 & 10.0 & 486.9 & 115.8 & 365.8 \\
 & exact, unbounded & 0.909 & 0.915 & 0.000 & 0.000 & 0.017 & 1.000 & 0.296 & 10.0 & 446.5 & 108.6 & 220.7 \\
 & widen & 0.907 & 0.913 & 0.000 & 0.000 & 0.017 & 1.000 & 0.306 & 9.89 & 772.1 & 594.4 & 1,572.7 \\
 & monitor 1\% & 0.896 & 0.900 & 0.003 & 0.017 & 0.017 & -- & 0.301 & 10.2 & 433.2 & -- & -- \\
 & tuple limit $10^6$ alone & 0.899 & 0.903 & 0.003 & 0.016 & 0.016 & -- & 0.299 & 10.0 & 429.4 & -- & -- \\
\midrule
MS MARCO, & report & 0.809 & 0.807 & 0.000 & 0.035 & 0.035 & -- & 0.356 & 6.07 & 10.7 & -- & -- \\
Qdrant & count & 0.809 & 0.807 & 0.000 & 0.035 & 0.035 & 1.000 & 0.356 & 6.02 & 26.5 & 21.7 & 28.2 \\
 & exact, $B$ 1,000 & 0.809 & 0.807 & 0.000 & 0.035 & 0.035 & 1.000 & 0.356 & 5.80 & 25.8 & 21.1 & 28.3 \\
 & exact, $B$ 10,000 & 0.809 & 0.807 & 0.000 & 0.035 & 0.035 & 1.000 & 0.356 & 5.95 & 26.4 & 21.6 & 27.9 \\
 & exact, unbounded & 0.833 & 0.842 & 0.000 & 0.000 & 0.035 & 1.000 & 0.343 & 6.10 & 44.4 & 37.9 & 50.4 \\
 & widen & 0.833 & 0.841 & 0.000 & 0.000 & 0.035 & 1.000 & 0.349 & 5.56 & 117.8 & 97.4 & 148.5 \\
 & monitor 1\% & 0.809 & 0.807 & 0.000 & 0.035 & 0.035 & -- & 0.356 & 5.15 & 10.6 & -- & -- \\
 & ACORN 1.0 alone & 0.987 & 0.986 & 0.000 & 0.000 & 0.000 & n/a & 0.085 & 15.9 & 125.7 & -- & -- \\
\bottomrule
\end{tabular}
\caption{X5: detecting and repairing under-retrieval (GitHub Actions run 40921769714; three builds per
store and corpus; every arm on every question). Recall@$k$ (R@5, R@10) and the rates are means over
retrievals with $k = 10$ unless marked. Empty and Short: the answer has no hit, or fewer than
$\min(k, R)$, where $R$ is what the subject may read. Trigger: the first search came back with fewer
than $k$ hits, which sets off the count. Detect: the share of truly short first searches flagged
\emph{missed}; n/a: there were none; ``--'': the arm does not count. FC: false-completeness rate, the
share of results marked complete that are not maximal. Latency: per-question median of three passes,
end to end; ``Triggered'' is the time spent on the short result (count and repair) on triggered
retrievals. Latency compares arms within one store and corpus only.}
\label{tab:rep-results}
\end{table*}

\paragraph{Detection (H5.1, confirmed).}
The count raised no false alarm and flagged every truly short first search as \emph{missed}: 992 of
992 for pgvector on the enterprise corpus, 574 of 574 for pgvector and 1,091 of 1,091 for Qdrant on MS
MARCO ($k = 5$ and $k = 10$ together), and every answer it marked \emph{exhausted} held all $\min(k,
R)$ readable chunks. Qdrant on the enterprise corpus had no truly short answer: its first searches came
back short only for people who could read fewer than $k$ chunks, which the count marked exhausted.
Qdrant did under-retrieve on MS MARCO, which X0 and X1 never saw on the enterprise corpus: 3.5\% of its
answers with $k = 10$ were short.

\paragraph{Exact repair (H5.2, confirmed for the unbounded arm).}
Repaired exactly without a bound, no answer was empty or short, and the repaired answers reached
recall@10 of 1.0000: on 576 triggered retrievals for pgvector on the enterprise corpus, on 323 for
pgvector and 644 for Qdrant on MS MARCO, which had returned 992, 574 and 1,091 short answers under
\emph{report} ($k = 5$ and $k = 10$ over three builds). Recall@10 rose from 0.667 to 0.782, from 0.900 to 0.915 and from 0.807 to 0.842. The
bound decided where repair applied. On the enterprise corpus every short answer had at most 1,000
readable chunks, so all three bounds repaired all of them; on MS MARCO a bound of 1,000 repaired 0.5\%
of pgvector's answers out of 1.7\% triggered, and every short answer of Qdrant had more than 10,000
readable chunks, so the bounded arms repaired none and reported each as \emph{over\_limit} and not
complete. At the default bound of 10,000, then, 0.1\% of pgvector's answers and all of Qdrant's 3.5\%
stayed short on MS MARCO. Where the repair ran without a bound, the count and the exact search
together took 30.1\,ms at the median
and 40.2\,ms at p95 of a triggered pgvector retrieval on the enterprise corpus, 108.6 and 220.7\,ms on
MS MARCO, and 37.9 and 50.4\,ms for Qdrant.

\paragraph{Cost on other queries (H5.3, inconclusive).}
Of the 60 median latency ratios between a repair arm and \emph{report} on the questions that
triggered neither, 45 lay within the pre-registered band of 0.95 to 1.05, one above it (1.054, Qdrant
on MS MARCO, build 2, $B$ 10,000) and 14 below it, down to 0.866: the repair arm faster than
\emph{report}. None exceeded the falsification bound of 1.10. Twelve of the 14 low ratios are
Qdrant's, and the arms ran in a fixed order after \emph{report} within each build, so they are
consistent with the store's latency drifting during a build rather than with any effect of the arms;
the run cannot separate the two. Over all questions, every repair arm's median latency lay within
0.6\,ms of \emph{report}'s.

\paragraph{Mis-retrieval (H5.4, confirmed).}
No repair touched Qdrant's loss on the enterprise corpus: 58.1\% of its answers marked complete under
\emph{report} were not maximal, and 57.3\% under every repair arm, a difference of $-0.008$ (95\%
interval $-0.027$ to 0.000). The small change is the count marking the short answers of people with
fewer than $k$ readable chunks exhausted, and so complete and correct; no answer was repaired. Across
the table, exact repair lowered the false-completeness rate only as far as short answers had caused it
(0.551 to 0.480 for pgvector on the enterprise corpus); the rest is mis-retrieval, which a full answer
hides from any count.

\paragraph{The monitor (H5.5, confirmed).}
The monitor sampled within the 99.9\% binomial interval around 1\% in all twelve builds (48 to 65 of
5,956 retrievals per build on the enterprise corpus, 210 to 267 of 24,800 on MS MARCO). Its recall,
computed against the store's own exact search, agreed within 0.1 with the harness's on all 1,767
sampled retrievals, and its mean estimated Qdrant's recall over all retrievals within 0.030 (0.684
on 166 samples against 0.654) on the enterprise corpus and within 0.004 (0.812 on 725 against 0.807)
on MS MARCO. A 1\% sample is thus enough to see the loss that \texttt{complete} cannot.

\paragraph{Widening (H5.6, confirmed).}
On the 421 questions that triggered in both arms, exact repair beat widening by 0.085 of recall@10
(95\% interval 0.064 to 0.110; Wilcoxon $p = 5.6 \times 10^{-30}$), and widening cost more: 360.9 and
1,572.7\,ms at p95 of a triggered pgvector retrieval, on the enterprise corpus and on MS MARCO, against
40.2 and 220.7\,ms for exact repair. A wider approximate search filled the short answers, since it
left none empty or short, but did not rank them exactly.

\paragraph{Store-native fixes (H5.7, confirmed).}
On the enterprise corpus Qdrant's per-tenant graphs reached recall@10 of 0.9386 while the best library
arm on the shared layout reached 0.6535, and the false-completeness rate fell from 0.581 to 0.235. ACORN
alone reached 0.855 there and 0.986 on MS MARCO, where it also removed every short answer, at a median
of 15.9 against 6.1\,ms. pgvector's tuple limit raised alone lowered the enterprise corpus's
empty-result rate only from 0.089 to 0.072, because the scan's memory bound then stops it, as X1 found
(Section~\ref{sec:qual-pgvector}).

The two tools divide the work. A count tells a short answer from a small readable set without error.
Exact repair without a bound fixes every short answer for tens to hundreds of milliseconds spent on
that answer; within a bound it fixes only the answers whose readable set fits, none of Qdrant's on MS
MARCO. Neither sees a full answer of the wrong chunks, which was most of the loss on both stores
(pgvector's recall@10 on the enterprise corpus was 0.782 after exact repair, 0.667 without it); only
a sampled exact search sees it, and only a better index or layout removes it.


\section{Side Channels in Filtered Retrieval}\label{sec:side}
% Side Channels in Filtered Retrieval: body of Section~\ref{sec:side} (the \section heading and its
% label stay in main.tex). Labels defined here: sec:side-game, sec:side-setup, tab:side-hyp,
% sec:side-results. Source: benchmarks/paper/PREREG_X3_X6.md on branch paper/x3-x6 (X6, with
% amendments 1 to 3, all made before any full result). Results marked \tbd{} come from the full run.

Proposition~\ref{prop:reduction} bounds what a subject can \emph{read}; it says nothing about what
the timing or shape of an answer reveals. B\"{u}ttcher and Clarke required of secure full-text
search that no aspect of the output, its order included, depend on content the user may not
search~\cite{buttcher2005security}. Each enforcement path of Section~\ref{sec:system} departs from
that requirement in a predictable way: a widening post-filter works longer when unreadable chunks
crowd the top of the ranking, an approximate filtered search passes unreadable nodes on its way to
readable ones, and a keyword score computed over the whole collection counts unreadable documents.
We ask how much an authenticated subject can learn from these departures about whether unreadable
chunks exist near a query, and what removing each one costs.

\subsection{A distinguishing game}\label{sec:side-game}

Two worlds of equal size differ only in where $m$ unreadable chunks lie. In $W_1$ they are the $m$
passages most similar to a probe query; in $W_0$ each is replaced by a passage of another domain, far
from every probe and cut to the same length, at the same position in the index. The attacker's
readable set is identical in both worlds. It issues the probe and sees only what a caller of the
layer sees: the latency, the number of hits, the completeness flag $c$, the mean score of the hits,
and the hits themselves, from which it can compute recall@10 against the exact top ten of its own
readable set, since it can read and embed that set. The \emph{differencing attacker}, our primary
adversary, sees the same probe in both worlds, as someone who repeats a query before and after a
document is added; its success is the paired AUC, the share of probes whose $W_1$ observation lies
further in the predicted direction than their $W_0$ observation, and its advantage is $2\,\mathrm{AUC}
- 1$. The \emph{one-shot attacker} sees one observation and must threshold it; its success is the AUC
over all $W_1$ against all $W_0$ observations.

\subsection{Setup}\label{sec:side-setup}

\paragraph{Data and probes.}
The background is all 57,638 FiQA passages~\cite{maia2018www18}, with its 6,648 queries as probe
candidates; far chunks come from 12,000 SciDocs abstracts, both from BEIR~\cite{thakur2021beir}, and
all text is embedded with all-MiniLM-L6-v2 (384 dimensions)~\cite{reimers2019sentencebert,wang2020minilm}.
Each of five seeds places eight groups of nine probes, one with $m = 1{,}000$ secrets, two with 100,
three with 10 and three with 1. A probe's secrets are its $m$ nearest passages, excluding hubs, the 1\% of
passages found among the 2,000 nearest neighbors of the largest number of queries; within a group, no probe's secrets fall in the
zone a group mate can explore, its nearest passages up to 1,000 beyond its own secrets, which is the
post-filter's cap. One index $W_1$ holds every group's secrets; for each group $g$, an index $W_0^g$
replaces only group $g$'s secrets by far chunks, so a probe of group $g$ compares two worlds that
differ in its own secrets alone. Background passages fall into 100 permission buckets, and three
attackers read 1\%, 10\% and 50\% of them; secrets carry an ACL none of them holds. Every seed
builds every index afresh.

\paragraph{Configurations.}
The \emph{pre-filter} runs each store's approximate filtered search at documented defaults:
pgvector, Qdrant and Elasticsearch, with FAISS over a flat index and \sys{}'s in-memory backend as
exact controls. The \emph{post-filter} is \sys{}'s widening loop at its defaults: fetch $4k$, widen
fourfold until $k$ hits are readable, the store runs out, or 1,000 candidates have been examined.
Four mitigations are measured against them. \emph{Constant-work widening} (M1) fetches the cap of
1,000 in one round, so every request examines the same number of candidates. \emph{Exact
pre-filtered search} (M2) ranks every readable chunk: pgvector with index scans disabled, Qdrant's
exact flag, Elasticsearch's \texttt{script\_score}. \emph{Separation} (M3) gives the attacker an index
that holds only its own tenant's chunks, as a tenant-per-index layout would; its advantage is zero by
construction, and measuring it twice provides an A/A null for the harness. \emph{Latency padding}
(M4) delays each answer to the next multiple of a budget fixed on the first seed and is evaluated on
the other four; it models an ideal padder and ignores scheduler jitter. For collection statistics,
keyword search runs under the pre-filter on Elasticsearch (BM25), on Qdrant (sparse vectors with an
IDF modifier) and on pgvector (\texttt{ts\_rank\_cd}, which uses no collection statistics: a
negative control), and Elasticsearch adds reciprocal-rank fusion of vector and keyword search.

\paragraph{Measurement and statistics.}
For every seed, store, probe, attacker and configuration, one warm-up is followed by five timed
repetitions, each of which runs $W_0^g$, $W_1$ and the two separated indexes in a fresh random order,
so drift affects every world alike; an observation is the median over repetitions. Every retrieval is
re-checked, and a hit outside the attacker's readable set would stop the run. Intervals are 95\%
cluster-bootstrap intervals over (seed, probe) pairs with 10,000 resamples. A channel is
\emph{present} when the paired AUC reaches the hypothesis's threshold and its interval's lower end
exceeds 0.55, and \emph{absent} when the interval lies within $[0.40, 0.60]$. Two guards precede any
verdict: the A/A comparison of the two separated indexes must show no latency channel, and exact
configurations must return identical hit lists in $W_0^g$ and $W_1$ for at least 99\% of probes. The
experiment runs on GitHub-hosted runners with four virtual CPUs and 16\,GB of memory.

\begin{table}[tbp]
\centering
\small
\begin{tabular}{@{}p{0.27\linewidth}p{0.47\linewidth}p{0.16\linewidth}@{}}
\toprule
Hypothesis & Prediction & Verdict \\
\midrule
H1 widening & Post-filter: count AUC $\ge 0.95$ at $m = 1{,}000$; latency AUC $\ge 0.75$ at $m = 100$
  ($p = 10\%, 50\%$); no channel at $m = 1$ & \tbd{} \\
H2 constant work (M1) & No latency channel; the count channel remains at $m = 1{,}000$ & \tbd{} \\
H3 approximate pre-filter & pgvector and Elasticsearch: latency AUC $\ge 0.75$ at $m = 1{,}000$;
  Qdrant: no channel at $p = 1\%, 10\%$ & \tbd{} \\
H4 exact pre-filter (M2) & No channel in any observable & \tbd{} \\
H5 collection statistics & BM25 and IDF-weighted sparse: score AUC $\ge 0.90$ at every $m$;
  \texttt{ts\_rank\_cd}: none; fusion changes the hit list for at least 25\% of probes at
  $m = 1{,}000$ & \tbd{} \\
H6 padding (M4) & No latency channel wherever the unpadded AUC was $\ge 0.75$ & \tbd{} \\
\bottomrule
\end{tabular}
\caption{Pre-registered hypotheses of the side-channel experiment, abbreviated; the full rules,
including what falsifies each, are in the released pre-registration.}
\label{tab:side-hyp}
\end{table}

\subsection{Results}\label{sec:side-results}

\tbd{X6 full run (GitHub Actions run 40940385704): guards G1 and G2; per hypothesis the verdict,
the paired and one-shot advantages per store, configuration, $m$ and $p$ with intervals; the
latency cost of M1, M2 and M4 and of separation; any leak (none expected). Figure: paired advantage
against $m$ per configuration.}


\section{End-to-End Evaluation With Canaries}\label{sec:e2e}
% End-to-End Evaluation With Canaries: body of Section~\ref{sec:e2e} (the \section heading and its
% label stay in main.tex). Labels defined here: sec:e2e-setup, sec:e2e-safety, tab:e2e-safety,
% sec:e2e-model. Source: evidence/kaggle-scale-run-2026-10-24.md (one run, wardowl 2.1.0, before the
% honest-completeness change). Repeats (X8) are pending and marked \tbd{}.

The previous sections test components in isolation. Here the whole pipeline runs as an organization
would run it: many tenants and people in one index, a policy composed of a tenant boundary and
per-document ACLs, adversarial questions, and a language model that answers from what retrieval
returns. The goal is a leak test with enough exposure to bound the leak rate, not a measure of
answer quality.

\subsection{Setup}\label{sec:e2e-setup}

\paragraph{Corpus.}
The index holds 107,540 chunks: 49,920 synthetic company documents in 20 tenants, of which 41,600
carry a unique canary string; 20 notices, one per tenant, each carrying a prompt injection; and
57,600 passages of the FiQA financial question-answering corpus~\cite{maia2018www18} from
BEIR~\cite{thakur2021beir}, which every tenant may read and which serve as distractors. Company
documents are drawn from templates (handbooks, finance, HR, board and project records) with
per-document ACLs over users and groups. Chunks are embedded with
bge-base-en-v1.5~\cite{xiao2024cpack}.

\paragraph{People and questions.}
189 people span the tenants. A typical employee may read the shared passages, the tenant's notice
and five of the twelve document templates, 58,641 chunks in all (the median); the least privileged
are non-executive board members (``judy'' personas), each of whom may read 624 chunks of three
templates in one tenant and none of the shared passages (0.58\% of the corpus); one person in no
group may read nothing. They ask 1,105 questions: 108 generic, 376 targeted at a document
they may read, 189 targeted at a document they may not read (\emph{forbidden}), and 432 red-team
prompts that ask for other people's data, other tenants' data or everything, or try to override
instructions.

\paragraph{Policy and stores.}
The policy is $\mathrm{AllOf}(\mathrm{Tenant}, \mathrm{MetadataACL})$ (Section~\ref{sec:sys-compose}).
It runs against ten configurations: FAISS over a flat index (exact, the reference), Chroma, pgvector,
pgvector with the same policy written in OPA, Qdrant, Elasticsearch, OpenSearch, Weaviate, Redis and
LanceDB, each at the settings listed in Appendix~\tbd{settings table}. A configuration with OpenFGA
behind the caching authorizer did not start (the engine's health endpoint did not answer within
180\,s) and is excluded. Each question retrieves with $k = 5$ for the answer and $k = 10$ for
recall@10. Qwen2.5-3B-Instruct~\cite{yang2024qwen25} answers from the five chunks.

\paragraph{Leak detection.}
A leak is any returned chunk outside the person's readable set, as decided by an oracle that
evaluates the policy without the layer, or any canary from an unreadable document found verbatim in
a retrieved chunk, in the prompt sent to the model or in the model's answer. Verbatim matching
misses a paraphrased secret; it is sufficient for retrieved chunks and prompts, which are copied, and
a lower bound for answers. \tbd{X8 adds paraphrase-aware matching.}

\subsection{Safety}\label{sec:e2e-safety}

\begin{table}[tbp]
\centering
\small
\begin{tabular}{@{}lrrrrr@{}}
\toprule
Configuration & Questions & Hits checked & Unreadable hits & Canaries (chunk/prompt/answer) & Violations \\
\midrule
FAISS (exact) & 1,105 & 16,360 & 0 & 0/0/0 & 0 \\
Chroma & 1,105 & 16,360 & 0 & 0/0/0 & 0 \\
pgvector & 1,105 & 16,178 & 0 & 0/0/0 & 0 \\
pgvector + OPA & 1,105 & 16,178 & 0 & 0/0/0 & 0 \\
Qdrant & 1,105 & 16,360 & 0 & 0/0/0 & 0 \\
Elasticsearch & 1,105 & 16,360 & 0 & 0/0/0 & 0 \\
OpenSearch & 1,105 & 16,360 & 0 & 0/0/0 & 0 \\
Weaviate & 1,105 & 16,360 & 0 & 0/0/0 & 0 \\
Redis & 1,105 & 16,360 & 0 & 0/0/0 & 0 \\
LanceDB & 1,105 & 16,360 & 0 & 0/0/0 & 0 \\
\midrule
Total & 11,050 & 163,236 & 0 & 0/0/0 & 0 \\
\bottomrule
\end{tabular}
\caption{End-to-end safety, one run. Hits are counted over both retrievals per question. pgvector
checked 182 fewer hits because some of its answers were short or empty (Section~\ref{sec:quality}).}
\label{tab:e2e-safety}
\end{table}

No configuration returned an unreadable chunk or exposed a canary (Table~\ref{tab:e2e-safety}). The
hits are not independent trials: they cluster by question and by person, and a fault in the layer
would affect a question, not a random hit. We therefore bound the leak rate per question, separately
for each configuration: with no leak in 1,105 questions, the one-sided 95\% upper bound is $1 -
0.05^{1/1105} \approx 2.7 \times 10^{-3}$, and per person $1 - 0.05^{1/189} \approx 1.6 \times
10^{-2}$ (Section~\ref{sec:prob-metrics}). Pooling the ten configurations would not shrink these
bounds honestly, since they share questions and the layer's code. The violation count is zero
because the filter and the check read the same ACL metadata here; Sections~\ref{sec:conformance}
and~\ref{sec:engines} test configurations in which they do not.

Two observations qualify the result. First, it is a statement about this layer under assumptions
(A1) to (A3) of Proposition~\ref{prop:reduction}, exercised with adversarial inputs, not a proof:
the red-team prompts can only change the query text, and the layer never lets text choose the
subject (Section~\ref{sec:sys-identity}). Second, safety held while quality did not. pgvector
returned no hits at all to ten least-privileged people who could read hundreds of chunks, and the
layer, before the change of Section~\ref{sec:sys-complete}, reported those empty answers as
complete; Qdrant answered, but with recall@10 of 0.56. Section~\ref{sec:quality} diagnoses both.
\tbd{X8: three repeated runs with fresh permission layouts; report the bound over all runs and
the run-to-run spread of recall.}

\subsection{What the model did with what it retrieved}\label{sec:e2e-model}

Prompt injection in readable content is outside our threat model (Section~\ref{sec:prob-threat}),
but the run measures where the boundary lies. Each tenant's injected notice is readable by that
tenant's employees and contractors, so retrieving it for them is correct. In the nine configurations
for which the run reports it, the notice was retrieved between 7 and 106 times, depending on the
store's ranking, and the 3B model followed its instruction every time it was retrieved (482 of
482). Retrieval-time enforcement decides which
content reaches the model, not what the model does with it; defenses against injected instructions
in permitted content~\cite{greshake2023not,debenedetti2025defeating} are complementary.

Answer quality followed retrieval quality. With FAISS the model found the target document of a
targeted question in its top five for every question and answered 99.7\% correctly; with Qdrant
the target was in the top five for 39.9\% of targeted questions. When a person asked about a
document they could not read, the model abstained in 70\% to 81\% of cases, depending on the store,
and otherwise answered from readable documents only, as the canary checks confirm.


\section{Related Work}\label{sec:related}
% Related Work: body of Section~\ref{sec:related} (the \section heading and its label stay in main.tex).
% The full account, with every precedent's figures and the industry systems, is in
% Appendix~\ref{app:related}. The enforcement design is not claimed as new; the contribution is
% verification and measurement.

Permission-aware retrieval draws on literatures that rarely cite one another; each settled part of the design we
test, and none verifies the whole. Appendix~\ref{app:related} gives the full account, including native store features
and commercial systems such as Amazon Bedrock's knowledge bases~\cite{aws2026access}, for which we found no independent
evaluation of permission correctness.

\paragraph{Enterprise search and databases.}
B\"{u}ttcher and Clarke gave the founding security model for multi-user full-text search: results ranked with global
statistics reveal unreadable files, and query integration fixes it~\cite{buttcher2005security}. Bailey et al.\ modeled late
binding and fetched growing blocks of candidates~\cite{bailey2006secure}. The pre-filter we test is query integration applied
to approximate search, and its widening post-filter is their loop; we claim neither. Query modification~\cite{stonebraker1974access},
Truman and non-Truman semantics~\cite{rizvi2004extending} and the sound, secure and maximal criteria~\cite{wang2007correctness} frame
what a pushed-down filter does and owes, and Kabra et al.\ showed that functions evaluated before the authorization predicate reveal
data~\cite{kabra2006redundancy}. Zerber~\cite{zerr2008zerber} and searchable encryption~\cite{song2000practical,curtmola2006searchable}
protect an index from an untrusted server, whereas we trust the store's host. This lineage proves rewriting correct under exact evaluation;
we test components that carry no such proofs, under approximate search.

\paragraph{Authorization.}
Zanzibar is late-bound and uses snapshot tokens against applying an old ACL to new content~\cite{pang2019zanzibar}; its descendants added
reverse lookups, with documented size limits~\cite{openfga2026search,authzed2026protecting}, and policy engines compile filters through
engineering artifacts~\cite{opa2026evaluating,cerbos2026cerbos}. Margrave~\cite{fisler2005verification}, Zelkova~\cite{backes2018semanticbased},
mutation testing~\cite{martin2007fault} and Cedar's differential testing~\cite{disselkoen2024how} address the policy and its check path. We found
no peer-reviewed semantics or correctness proof for any policy-to-filter translator, and we measure, at one revision, how often filters disagree
with their checks.

\paragraph{Filtered search and permission-aware indexes.}
The collapse of post-filtering is well measured~\cite{patel2024acorn,li2025attribute,kim2026racorn1}. Index designs for attribute-filtered
search~\cite{wang2023efficient,gollapudi2023filtereddiskann,cai2024navigating,li2025sieve} take the predicate as part of the query rather than
compile it from an external policy; benchmarks sweep selectivity and control correlation~\cite{zhang2026vecbench,amanbayev2026filtered}, and HCBGen
shows that neither alone is a stable proxy for difficulty~\cite{lim2026revisiting}. None uses real ACLs. Permission-specific indexes are
recent~\cite{zhong2026honeybee,han2026dont}, and a vision paper asks for correctness criteria~\cite{yalamarthi2026policyaware}. We propose no index and
claim neither the first permission-aware vector search nor alignment as a new axis. We found no evaluation of per-query disjunctions over hundreds or
thousands of group labels, and no audit of production stores under such filters at documented, tuned and exact settings.

\paragraph{RAG privacy and agents.}
Anyone who can query a RAG system can drain its store~\cite{zeng2024good,qi2024follow}, and model-side defenses are
probabilistic~\cite{bhatt2025enterprise}. Attribute-based gating at retrieval removed cross-tenant leakage~\cite{arceo2026securing}, and re-checking
at the graph-expansion boundary cut a retrieval-pivot risk~\cite{thornton2026retrieval}. Corpus poisoning~\cite{zou2025poisonedrag} and information-flow
control for agents~\cite{costa2025securing} are complementary, and token-bound identity~\cite{jones2020oauth,mcp2026model} is adopted here, not claimed.
This literature shows why enforcement must be deterministic and precede generation, but does not test whether a layer is faithful or what it costs in
recall.

\paragraph{Side channels and testing.}
Our game tests a restricted form of noninterference~\cite{goguen1982security} over chosen observables. Timing has long revealed hidden content on the
web~\cite{felten2000timing,gelernter2015crosssite}, and padding is the simplest case of predictive mitigation~\cite{askarov2010predictive}; we add a measurement
of how much each enforcement path of filtered vector search reveals. Database engines have mature logic-bug
oracles~\cite{rigger2020testing,rigger2020detecting,rigger2020finding}, while vector databases have none~\cite{xie2025toward,wang2025towards}. Our testing is
differential in McKeeman's sense~\cite{mckeeman1998differential}, and it evaluates its relations under exact search, so that approximation cannot confound them.


\section{Limitations, Ethics and Disclosure}\label{sec:limits}
% Limitations, Ethics and Disclosure: body of Section~\ref{sec:limits} (the \section heading and its
% label stay in main.tex). Labels defined here: sec:limits-limits, sec:limits-ethics, sec:limits-ai,
% sec:limits-repro. Items marked \tbd{} must be settled before the preprint.

\subsection{Limitations}\label{sec:limits-limits}

\paragraph{Synthetic permissions.}
No public corpus pairs documents with the access-control lists of a real organization, so every
permission layout here is generated: tenants, groups and per-document ACLs drawn from templates and
seeded generators, over public passages (FiQA) and templated company documents. Real
organizations may have more skewed layouts, with many small groups, deep group hierarchies and
frequent exceptions. Two choices limit the damage. We place every (question, subject) pair on
measured correlation axes instead of trusting a layout's label (Section~\ref{sec:prob-metrics}), so
results transfer to any layout with the same measured difficulty; and we vary layouts across seeds.
The enterprise scenario is harsh in one respect that we state: each company document has 19
near-identical copies in other tenants that only the tenant rule separates.

\paragraph{One instrument.}
Every measurement passes through one enforcement layer, which the authors wrote, so a fault in its
filter translators could be mistaken for a store fault. The conformance kit compares every answer
with an in-memory reference that shares no code with the translators, and the retrieval audit uses
each store's own exact search as a second baseline. Of the faults we attribute to a store or engine,
those found by our experiments come with a reproducer that uses only the vendor's own client, and
each showed the reported behavior when run on GitHub-hosted runners on 30 October 2026 against the
version we evaluated and the vendor's latest release (Redis 8.10.2 and 8.10.3, MongoDB Atlas Local
8.0.30 and 8.3.13, Cerbos 0.56.0 and 0.57.0, and Chroma 1.5.9, still its latest). Running them
corrected one figure: MongoDB's vector-search filter fields keep 8,191 characters, not the 8,181 we
had taken from its documentation of another field type. The other rows rest on the vendors'
documentation, which Table~\ref{tab:conf-faults} cites.

\paragraph{Versions, settings and hardware.}
Results hold for the versions and settings we state (for example pgvector 0.8.7 and Qdrant 1.19.2)
and may not survive the next release; vector stores change quickly, and several behaviors we
measure are documented defaults that vendors may revise. Experiments ran on shared cloud hardware
(Kaggle sessions with two T4 GPUs and four CPUs, and GitHub-hosted runners), so latencies are
comparable within an experiment, not across experiments or with vendor benchmarks.

\paragraph{Statistical power.}
The end-to-end scenario ran four times with fresh documents and questions, and four more times for
the OpenFGA configuration, but always with one permission layout; its people recur in every run, so
its bound per person is a single run's. The diagnosis grid ran three index builds per layout, the
frontier and repair experiments three builds per store and corpus over ten permission layouts,
and the engine, filter-size and side-channel experiments five seeds each; intervals come from
bootstrap resampling over people, subjects or probes, not over builds, so a store whose recall varies
between builds (Qdrant's by up to 0.086) is described by the mean of three. Zero-leak results are
reported as upper bounds per question and per subject, not per hit. Several pre-registered verdicts
are inconclusive rather than confirmed or falsified, and we report them as such.

\paragraph{Scope of the security claim.}
Proposition~\ref{prop:reduction} needs correct checks, faithful metadata and failing closed. A store
that binds content to the wrong identifier or metadata defeats any re-check, and we test that
assumption only as far as the conformance kit's metadata checks reach. The completeness flag
detects under-retrieval but not mis-retrieval: a full answer with the wrong chunks still reports
complete. Side channels are measured (Section~\ref{sec:side}), not eliminated.

\subsection{Ethics and responsible disclosure}\label{sec:limits-ethics}

No human subjects took part. People, documents, ACLs and canaries are synthetic, and the shared
passages come from a public research corpus~\cite{maia2018www18}. Red-team prompts are generic
attempts to obtain other people's data and target only our own deployments; every store and engine
ran in containers or processes we controlled, on our own sessions and runners, and no third-party
service was probed.

None of our findings is a leak through the layer, because the re-check contains filter faults. They
can still matter for deployments that trust a pushed-down filter without a re-check: a filter that
admits more than the check would leak there, and one, Chroma's crash on a large filter, lets any
user in enough groups stop an application that passes their groups to Chroma. For each fault in a
third-party component that our experiments found, we prepared a report with a minimal reproducer
that does not need our library, and Table~\ref{tab:conf-faults} gives each report's status at the
time of writing. \tbd{filing dates, links and vendor responses, once the reports are filed}

\subsection{Use of AI assistance}\label{sec:limits-ai}

The library, the conformance kit and the experiment harnesses were written largely by an AI coding
agent under the authors' direction, and the agent also assisted with literature search and with
drafting this text. The authors chose the questions, pre-registered the predictions, reviewed the
designs and are responsible for every claim. Each claim rests on executed tests or logged runs whose
outputs we release, and each reference was checked against a primary record (title, authors, venue
and year), with unverifiable fields omitted rather than guessed. \tbd{Final human check of every
reference before submission.} AI tools are not authors. This statement follows the venue's policy on
large language models and arXiv's policy on generative AI.

\subsection{Reproducibility}\label{sec:limits-repro}

We release the library at the evaluated version, the conformance kit, every harness script, the
Kaggle notebooks and GitHub workflows that ran them, the pre-registrations with their dated
amendments, and the raw outputs with container image digests and seeds \tbd{archive DOI}. For
review, an anonymized mirror with the package renamed is provided. Compute: about 35 hours of
Kaggle sessions (two T4 GPUs or four CPUs each) and more than 140 hours of GitHub-hosted runner jobs,
listed per experiment in Appendix~\ref{app:compute}; no experiment reported here ran on a local
machine.


\section{Conclusion}\label{sec:conclusion}
% Conclusion: body of Section~\ref{sec:conclusion}. Every number is from the sections it summarizes.

Permission-aware retrieval is usually described as a filter. We have argued that it is better
treated as a composition of untrusted parts, translators, approximate indexes and filter APIs, around
a small trusted computing base centered on the check API. Re-checking every hit with that API makes
containment independent of the untrusted parts: under the stated assumptions no returned chunk lies
outside the readable set at the revision the check read. Our measurements agree, with no leak in
any of eleven end-to-end configurations, including red-team prompts (one-sided 95\% bounds of
$6.8 \times 10^{-4}$ per question and $1.6 \times 10^{-2}$ per person, over the question generator's
distribution), and none in 1,050 fault episodes ($2.8 \times 10^{-3}$ per episode, over the faults we
injected).

Containment, however, is the easy part, and it is not all that security asks. The revision can be
stale: revoked permissions kept admitting for up to a cache's time to live or an engine's staleness,
and an invalidation in front of a stale engine broke the bound the release stated until we fixed it.
Answers still reveal: latency, counts, recall and scores told an attacker when unreadable chunks lay
near a query. And the untrusted parts still decide how much of what a person may read they actually
receive, where the stack is far from faithful. The layer's translation of one engine's query plans
admitted what its checks denied on data of an unexpected shape, stores dropped matching chunks at size
limits and one crashed its client on a long filter, and filtered approximate search lost much of a
readable set when that set was far from the question: on clustered layouts with negatively correlated
questions at their defaults, five of nine stores answered filters admitting 0.5\% to 10\% of a million
passages worse than the smallest, and the same pgvector configuration was exact in some runs and lost
about a quarter of its recall in others, a split consistent with index maintenance that no setting
names. In the release we ran end to end, none of this reached the application as a signal. It can be
made to: a filter-only count flagged every truly short answer without a false alarm, and an exact
search without a bound repaired them all, while a bound leaves unrepaired those above it, as with
Qdrant on MS MARCO; but no count sees a full answer of the wrong chunks, which was most of the remaining
loss, and only a sampled exact search or a better layout addresses it.

Three practices follow. Test filter semantics differentially, against a reference that shares no code
with the translators, and measure recall under permission-shaped filters against exact search over
the readable set, on measured correlation axes, before choosing a store, a layout or a plan. Never let
a system call a short approximate answer complete: count, and fall back to exact search within a
bound, or say that results may be missing; and sample exact searches to see what a count cannot.
Treat the shape and timing of an answer as outputs too, and the staleness of a check as part of the
claim: state $\delta$ for each path, and where unreadable chunks near a query must not be inferable,
search exactly and pad latency. We release the kit, the harnesses and the data so that these checks
can be repeated as stores and engines change.


\bibliographystyle{plainnat}
\bibliography{../bib/references,../bib/web}

\appendix
% Appendices. Labels defined here: app:settings, tab:app-x0-settings, app:x1grid, app:kit, app:compute.
% Sources: evidence/kaggle-scale-run-2026-10-24.md, investigations/filtered-recall-gaps.md (sections
% 1.3, 1.4 with file:line references), evidence/conformance-campaign-run1.md,
% benchmarks/paper/PREREGISTRATION.md on the X1 branch. Versions marked \tbd{} are read from the
% archived outputs before the preprint.

\section{Store settings in the enterprise scenario}\label{app:settings}

\begin{table}[h]
\centering
\scriptsize
\setlength{\tabcolsep}{4pt}
\begin{tabular}{@{}p{0.17\linewidth}p{0.2\linewidth}p{0.55\linewidth}@{}}
\toprule
Store & Version & Index and search settings \\
\midrule
FAISS & faiss-cpu 1.15.1 &\texttt{IndexFlatIP} (exact); metadata matched in Python \\
pgvector & pgvector 0.8.7, PostgreSQL 16.15 & HNSW built after loading with default $m = 16$,
  \texttt{ef\_construction} 64; \texttt{hnsw.ef\_search} 100; iterative scan \texttt{relaxed\_order};
  \texttt{hnsw.max\_scan\_tuples} 20,000 and \texttt{scan\_mem\_multiplier} 1 (defaults);
  \texttt{work\_mem} 4\,MB; metadata in \texttt{jsonb} with a GIN \texttt{jsonb\_path\_ops} index;
  filters as \texttt{@>} containment tests \\
Qdrant & 1.19.2 & HNSW cosine, $m = 16$, \texttt{ef\_construct} 100, full-scan threshold 10,000\,KB
  (defaults); keyword payload indexes on users, groups, tenant and identifier, created before
  loading, none marked \texttt{is\_tenant}; \texttt{hnsw\_ef} 128; no ACORN, not exact \\
Elasticsearch & 9.5.5 & \texttt{int8\_hnsw}, \texttt{num\_candidates} 500 \\
OpenSearch & 3.9.0 & Lucene HNSW, \texttt{num\_candidates} 500 \\
Weaviate & 1.39.12 & HNSW (defaults), indexed while importing \\
Redis & 8.10.3 & HNSW, \texttt{EF\_RUNTIME} 100 \\
LanceDB & lancedb 0.41.0 & \texttt{IVF\_PQ}, \texttt{nprobes} 50, \texttt{refine\_factor} 30 \\
Chroma & 1.5.9 & HNSW cosine (defaults) \\
\bottomrule
\end{tabular}
\caption{Settings in the end-to-end runs (Sections~\ref{sec:e2e} and~\ref{sec:quality}). Each run is
one Kaggle session, two T4 GPUs, four CPUs, 31.3\,GiB of memory; embeddings bge-base-en-v1.5, 768
dimensions, unit length. The setup scripts pin the server versions of Qdrant, Elasticsearch,
OpenSearch, Weaviate, PostgreSQL's pgvector, OPA and OpenFGA; Redis comes from its package repository
and the embedded stores from PyPI, at the versions the three repeat runs recorded.}
\label{tab:app-x0-settings}
\end{table}

\paragraph{Ground truth.}
The exact baseline multiplies all question vectors with all chunk vectors, masks chunks the oracle
says the person may not read, and takes the top ten. FAISS over a flat index reproduces it at recall
1.000, which validates the pipeline. A hit counts toward recall if it is readable and its similarity,
recomputed from the run's vectors, is at least the $k$-th exact score minus $10^{-4}$. Recall@5 and
recall@10 come from separate retrievals with $k = 5$ and $k = 10$. People who may read nothing are
excluded.

\section{The diagnosis grid}\label{app:x1grid}

\paragraph{pgvector.} The numbers in the paper are from run 2: three index builds with HNSW $m = 16$
and \texttt{ef\_construction} 64 and three with $m = 32$ and \texttt{ef\_construction} 200 (a core
grid), 1,489 questions (the scenario's 1,105 and 384 asked by 24 people in a single group each),
three timed passes each. Query-time settings: the scenario's settings with the planner's choice and
prepared statements (\emph{as run}); custom plans forced (\texttt{plan\_cache\_mode =
force\_custom\_plan}); no prepared statements; the HNSW scan forced (custom plans with sequential
and bitmap scans disabled), the reference for the scan-budget hypotheses, and every budget and depth
setting forced the same way: \texttt{hnsw.max\_scan\_tuples} of 20,000, 100,000 or 1,000,000 crossed
with \texttt{hnsw.scan\_mem\_multiplier} of 1, 8 or 16; \texttt{work\_mem} 64\,MB;
\texttt{hnsw.ef\_search} 40, 200, 400 and 1,000; an inner \texttt{LIMIT} of four or ten times $k$;
strict ordering; and exact search with index scans disabled. Two replays of the audit's pass ran
first on the first build, right after loading and after \texttt{VACUUM (ANALYZE)}; every other
setting ran after a vacuum. PostgreSQL's \texttt{auto\_explain} recorded the plan of every search of
the first pass. Run 1 built one pgvector index, skipped 12 settings for time, and recorded no
executed plans; it is reported in the released evidence.

\paragraph{Qdrant.} Five layouts (as run; tenant field marked \texttt{is\_tenant}; per-tenant graphs
with \texttt{payload\_m} 16 and $m = 0$; shared passages in a separate collection; one segment), three
builds each, with \texttt{hnsw\_ef} 64, 128 and 512, ACORN with maximum selectivity 0.4 and 1.0, and
exact search. Search paths per segment were read from Qdrant's telemetry.

\paragraph{Statistics.} 95\% intervals from 10,000 cluster-bootstrap resamples over people (seed
20261026); paired differences between settings on the same questions; verdicts computed by script
from the pre-registered criteria.

\section{The conformance kit's data variants}\label{app:kit}

Each store target ran 20 seeds of nine variants: the kit's default data (ten corpora per seed) and
eight stress variants (three corpora per seed): corpora of up to 150 documents; filters nested six
deep; disjunctions over 100 and 1,000 values; strings up to 40,000 characters (long and huge);
case and Unicode confusables; numeric strings next to integers; and unusual field names. Each engine
target ran 20 seeds of seven variants: default; case and Unicode; numeric strings; 1,000 groups;
300 chunks per check; long identifiers; and scalar instead of list ACL fields. A run is one (seed,
variant) and takes its worst subject's outcome: pass, a conformance violation, a refusal by the
library (a \texttt{ValueError} for data it cannot represent), or an error. The in-memory reference
store and the in-process metadata-ACL engine are controls.

\section{Compute}\label{app:compute}

\begin{table}[h]
\centering
\small
\begin{tabular}{@{}llrr@{}}
\toprule
Experiment & Platform & Wall clock & Session or job hours \\
\midrule
End-to-end, first run (seed 7) & Kaggle, $2\times$T4, 4 CPUs & 2.8\,h & 2.8 \\
End-to-end, three repeats (X8) & Kaggle, $2\times$T4, 4 CPUs, one session each & 2.8\,h each & 8.4 \\
End-to-end, OpenFGA, four runs & Kaggle, $2\times$T4, 4 CPUs, one session each & 0.9\,h each & 3.6 \\
Diagnosis grid, run 1 (X1) & Kaggle, $2\times$T4, 4 CPUs & 7.3\,h & 7.3 \\
Diagnosis grid, run 2 (X1) & Kaggle: 1 GPU session, 4 CPU sessions & 6.1\,h & 12.5 \\
Conformance campaign, run 1 & GitHub-hosted runners, 6 jobs & 2.9\,h & n/a \\
Conformance campaign, run 2 & GitHub-hosted runners, 10 jobs & 2.1\,h & 8.2 \\
Engines (X4, X7) & GitHub-hosted runners, 15 jobs & 2.0\,h & 15.3 \\
Revocation and faults, rerun (X7) & GitHub-hosted runners, 8 jobs & 1.5\,h & 1.5 \\
Filter size and side channels (X3, X6) & GitHub-hosted runners, 15 jobs & 3.3\,h & 10.2 \\
Their follow-up & GitHub-hosted runners, 11 jobs & 2.2\,h & 7.4 \\
Frontiers and repair (X2, X5) & GitHub-hosted runners, 72 jobs & 10.6\,h & 98.7 \\
\bottomrule
\end{tabular}
\caption{Compute of the runs reported in the paper. Kaggle sessions have 4 CPUs and 31.3\,GiB of
memory; GitHub-hosted runners have 4 virtual CPUs and 16\,GB. Job hours sum the jobs' run times; wall
clock includes queueing. Smoke runs that tested the harnesses, whose numbers were never analysed,
and runs cancelled before they produced results are not listed.}
\label{tab:app-compute}
\end{table}


\end{document}
